\documentclass[11pt,reqno]{amsart}

\usepackage{amsmath,amssymb,amsthm}
\usepackage{mathtools}
\usepackage[margin=1in]{geometry}
\usepackage[expansion=false]{microtype}
\usepackage[colorlinks=true,linkcolor=blue,citecolor=blue]{hyperref}

%% ---- theorem environments ----
\newtheorem{theorem}{Theorem}[section]
\newtheorem{proposition}[theorem]{Proposition}
\newtheorem{lemma}[theorem]{Lemma}
\newtheorem{corollary}[theorem]{Corollary}
\theoremstyle{definition}
\newtheorem{definition}[theorem]{Definition}
\newtheorem{conjecture}[theorem]{Conjecture}
\newtheorem{remark}[theorem]{Remark}
\newtheorem{assumption}[theorem]{Hypothesis}

%% ---- macros ----
\newcommand{\W}{\mathcal{W}}
\newcommand{\Jc}{\mathcal{J}}
\newcommand{\Cor}{\mathcal{C}}
\newcommand{\spec}{\operatorname{spec}}
\newcommand{\specess}{\operatorname{spec_{ess}}}
\newcommand{\negind}{\operatorname{neg.ind}}
\newcommand{\ran}{\operatorname{ran}}
\newcommand{\dist}{\operatorname{dist}}
\newcommand{\tr}{\operatorname{tr}}
\newcommand{\rank}{\operatorname{rank}}
\newcommand{\R}{\mathbb{R}}
\newcommand{\C}{\mathbb{C}}
\newcommand{\N}{\mathbb{N}}
\newcommand{\Z}{\mathbb{Z}}
\newcommand{\Hil}{H}
\newcommand{\sst}{\operatorname{st}}
\newcommand{\Qaut}{Q^{\mathrm{aut}}}
\newcommand{\half}{\tfrac12}
\DeclareMathOperator{\sgn}{sgn}

\begin{document}

\title[The Weil--Toeplitz algebra of window forms]{The Weil--Toeplitz algebra of window forms and a logistic tensor-tower identity}

\author{David Alejandro Trejo Pizzo}
\thanks{Independent Researcher. \texttt{dtrejopizzo@gmail.com}}

\begin{abstract}
We organize renormalized window forms attached to the Weil explicit formula as a
bounded field of finite-dimensional Hilbert spaces over \(\N\); the bounded
sections form a \(C^*\)-algebra \(\W\) (an \(\ell^\infty\)-product), the
sections vanishing at infinity form an ideal \(\Jc\), and the class of a
section in the corona \(\Cor=\W/\Jc\) is its \emph{symbol}, whose spectrum is
the limit set of the window spectra. Within an axiomatized spectral model
\(\mathfrak W\) we prove a trichotomy: off-critical quadruples of the model are
seen by the symbol (fat), accumulate at the point \(0\) of its spectrum (thin
visible), or leave the symbol unchanged (sub-resolution); we also record that
strict essential positivity is exactly the form of positivity that survives
compact perturbations. Independently we prove the exact identity
\(\sigma^{(k)}=\half(1-(1-2\sigma^0)^{2^k})\) for normalized negative indices of
tensor powers, and apply it to a localized quantity \(\sigma_{\mathrm{loc}}\)
built from a doubled form \(Q_2\) defined in Section~5 by a spectral sum over
pairs of zeros: \(\sigma_{\mathrm{loc}}=0\) under RH, and
\(\sigma_{\mathrm{loc}}>0\) under \(\neg\)RH, using an explicit family of
B-spline-type bumps and a localization radius above an absolute constant. Since
\(Q_2\) is defined through the zeros, the resulting equivalence of RH with
\(\sigma_{\mathrm{loc}}=0\) is close to a restatement of the existence of an
off-critical zero. The paper does not prove RH, and the connection of \(Q_2\)
and of the window model with the primes is not established here.
\end{abstract}

\maketitle

\tableofcontents
\newpage

\section{Introduction}\label{sec:intro}

\subsection{The framework in one sentence}
The Weil explicit formula defines, on a class of test functions equipped with a
notion of height, a hermitian form \(Q\) whose positivity is equivalent to the
Riemann Hypothesis (RH) \cite{Weil1952}. We package the renormalized
\emph{windows} of this form as a section of a bounded field of
finite-dimensional Hilbert spaces over \(\N\), form the associated
\(C^*\)-algebra \(\W\) of
bounded sections together with its corona \(\Cor=\W/\Jc\) (Definition
\ref{def:field}), and read the obstruction to RH inside the spectrum of the
\emph{symbol} \(\sigma(Q)\in\Cor\). This is the exact analogue, for a family,
of the passage from a Toeplitz operator to its Calkin image: the essential
spectrum is computed window-by-window in the limit.

\subsection{Content}
There are two independent parts.

\emph{(A) The window algebra.} The corona \(\Cor\) turns the gap between what
is controlled in mean over windows and what is needed uniformly over windows
into membership in \(c_0\) versus \(\ell^\infty\), equivalently into the
difference between the Ces\`aro states of \(\Cor\) and its norm (Corollary
\ref{cor:master-quantifier}, Section \ref{sec:honest}(b)). Inside an axiomatized model
we prove a trichotomy (Theorems \ref{thm:fat}--\ref{thm:nogo}): off-critical
quadruples of the model stratify by visibility, and the intermediate stratum
lives at the boundary point \(0\in\spec\sigma(Q)\). By Theorem \ref{thm:exact}
(a standard characterization, proved here for completeness) this is the point
separating the positivity that survives compact perturbations from the one that
does not.

\emph{(B) A logistic identity and a localized quantity.} The normalized negative
indices of the tensor powers of a nondegenerate hermitian form obey an exact
logistic recursion (Theorem \ref{thm:logistic}). Applied to a localized quantity
\(\sigma_{\mathrm{loc}}\) built from a doubled form \(Q_2\) (defined in Section
\ref{sec:logistic} by the spectral sum \eqref{eq:speca} over pairs of zeros), it
gives \(\lim_k\sigma_{\mathrm{loc}}^{(k)}=0\) under RH (Proposition
\ref{prop:RHbranch}), and \(\lim_k\sigma_{\mathrm{loc}}^{(k)}=\half\) under
\(\neg\)RH (Theorem \ref{thm:realiz}). The \(\neg\)RH branch uses an explicit
family of bumps (Lemma \ref{hyp:pw}) and a localization radius above an
absolute constant; radius \(2\) does not suffice for the tail estimate with any
family (Remark \ref{rem:pw-rigidity}). Because \(Q_2\) is defined by a sum over
the zeros, the resulting equivalence is a localized, finite-dimensional
restatement of the existence of an off-critical zero, not an arithmetic
criterion.

\subsection{Scope}\label{subsec:honest-up-front}
(a) We do not prove RH. The equivalence of Theorem \ref{thm:main} between RH and
\(\sigma_{\mathrm{loc}}=0\) is proved for the zero-side form \(Q_2\) of this
paper; since \(Q_2\) is defined through the zeros, it says that a configuration
with an off-critical zero is detected by some cell specified without zero
positions, and it carries no arithmetic information. (b) Whether the symbol \(\sigma(Q)\) of the genuine
Weil windows equals an arithmetic prediction is equivalent to a uniform
second-order asymptotic over windows (Theorem \ref{thm:symbolcompute}); it is
not resolved. (c) The model axioms (S1)--(S3) and Hypotheses
\ref{hyp:blocks}--\ref{hyp:autonomy} are not proved for \(\zeta\), and an
expression of \(Q_2\) through primes is not established. Results labelled
Theorem, Proposition or Lemma are proved here (some are standard and attributed);
conjectures and hypotheses are stated as such.

\subsection{Relation to earlier work}
The windowed hypotheses below are motivated by the structure of the localized
Weil form on even test functions in \cite{P18}: each off-critical quartet
contributes one Lorentzian block of signature \((1,1)\) \cite[Theorem~3.1]{P18},
and the negative index is at most the number of quartets
\cite[Corollary~6.5]{P18}. The windowed statements used here are not proved in
\cite{P18}; they are stated as hypotheses.

\subsection*{Notation}
\(H\) denotes a complex separable Hilbert space; \(\mathcal B(H)\) the bounded
operators; \(G=G^*\) a bounded self-adjoint operator with spectral measure
\(E_G\). A hermitian form is sesquilinear, linear in the first variable, with
\(Q(y,x)=\overline{Q(x,y)}\). For a subspace \(W\), \(\negind(Q|_W)\) is the
supremum of \(\dim N\) over subspaces \(N\subseteq W\) with \(Q(x,x)<0\) for all
\(x\in N\setminus\{0\}\), valued in \(\N\cup\{\infty\}\). A nontrivial zero is
written \(\rho=\beta+i\gamma\), \(\beta\in(0,1)\); a zero is \emph{off-critical}
if \(\beta\neq\half\); a quadruple is \(\{\half\pm\delta\pm i\gamma\}\).

\section{Defect spaces and the defect \texorpdfstring{$\delta$}{delta}}
\label{sec:defect}

We isolate the minimal language of index needed below. The defect of a hermitian
form is the dimension of a maximal negative subspace; it is the
polarization-theoretic analogue of the Fredholm index.

\begin{definition}[defect]\label{def:defect}
For a hermitian form \(Q\) on a complex vector space \(H\) with a marked
subspace \(P\), the \emph{defect} is
\(\delta(Q,P):=\negind(Q|_P)\in\N\cup\{\infty\}\). When \(H\) is Hilbert,
\(P=H\), and \(Q(x,y)=\langle Gx,y\rangle\) with \(G=G^*\) bounded, we write
\(\delta(G)\). The form is \emph{polarized} if \(\delta=0\).
\end{definition}

\begin{lemma}[the defect is spectral]\label{lem:defect-spectral}
\(\delta(G)=\dim\ran E_G(-\infty,0)\).
\end{lemma}

\begin{proof}
On \(N_-:=\ran E_G(-\infty,0)\) one has
\(\langle Gx,x\rangle=\int_{(-\infty,0)}\lambda\,d\|E_\lambda x\|^2<0\) for
\(x\neq0\); thus \(N_-\) is negative and \(\delta\geq\dim N_-\). Conversely, if
\(N\) is negative then \(E_G(-\infty,0)|_N\) is injective: if
\(E_G(-\infty,0)x=0\) then \(x\in\ran E_G[0,\infty)\), whence
\(\langle Gx,x\rangle\geq0\), impossible for \(x\in N\setminus\{0\}\). So
\(\dim N\leq\dim N_-\).
\end{proof}

\subsection{Arithmetic of the defect}
The defect is monotone, additive, and tensor-multiplicative in a precise sense.
These are used only through the \emph{tensor identity}, which drives the
logistic law of Section \ref{sec:logistic}; we re-prove what we use.

\begin{lemma}[monotonicity and additivity]\label{lem:mono-add}
\textup{(i)} If \(f\colon(H_1,Q_1,P_1)\to(H_2,Q_2,P_2)\) is linear, isometric
\((Q_2(fx,fy)=Q_1(x,y))\) with \(f(P_1)\subseteq P_2\), then
\(\delta(Q_1,P_1)\leq\delta(Q_2,P_2)\). \textup{(ii)}
\(\delta(Q_1\oplus Q_2)=\delta(Q_1)+\delta(Q_2)\), with no dimension or
nondegeneracy hypothesis.
\end{lemma}

\begin{proof}
(i) For \(N\subseteq P_1\) negative, \(f|_N\) is injective (\(fx=0\Rightarrow
Q_1(x,x)=Q_2(fx,fx)=0\Rightarrow x=0\)) and \(f(N)\subseteq P_2\) is negative by
isometry. (ii) ``\(\geq\)'' is direct sum of witnesses. ``\(\leq\)'': if some
\(\delta(Q_i)=\infty\) there is nothing to prove; if both are finite, choose
negative subspaces \(N_i\subseteq P_i\) of maximal dimension and let
\(R_i:=\{w\in P_i:Q_i(w,x)=0\ \forall x\in N_i\}\); since \(Q_i\) is definite on
the finite-dimensional \(N_i\), \(P_i=N_i\oplus R_i\). Then \(Q_i\geq0\) on
\(R_i\): if \(w\in R_i\) had \(Q_i(w,w)<0\), then by orthogonality
\(N_i\oplus\C w\) would be negative of larger dimension. For \(N\subseteq P_1\oplus P_2\) negative, the projection
\(\pi\colon N\to N_1\oplus N_2\) along \(R_1\oplus R_2\) has trivial kernel
(\(\ker\pi=N\cap(R_1\oplus R_2)\) is both \(\geq0\) and \(<0\)), so
\(\dim N\leq\delta(Q_1)+\delta(Q_2)\).
\end{proof}

\begin{lemma}[Whitney formula for the defect]\label{lem:whitney}
For finite-dimensional nondegenerate hermitian forms of signatures
\((p_1,q_1)\), \((p_2,q_2)\) (with \(q_i=\delta\)),
\[
\sgn(Q_1\otimes Q_2)=(p_1p_2+q_1q_2,\;p_1q_2+q_1p_2),
\qquad
\delta(Q_1\otimes Q_2)=p_1q_2+q_1p_2 .
\]
\end{lemma}

\begin{proof}
Take \(Q_i\)-orthonormal bases with signs \(\varepsilon^{(i)}_a\in\{\pm1\}\)
(Sylvester's law of inertia \cite{Gantmacher}). The tensor basis is \(Q_1\otimes Q_2\)-orthogonal with signs
\(\varepsilon^{(1)}_a\varepsilon^{(2)}_b\); the sign is \(-1\) exactly when the
two factors disagree, giving \(p_1q_2+q_1p_2\) negative directions and a
diagonal orthogonal basis with no isotropic vectors (nondegeneracy).
\end{proof}

\begin{corollary}[reduced signature is the only multiplicative character]
\label{cor:character}
With \(d=p+q\), \(s=p-q\): \(d(Q_1\otimes Q_2)=d_1d_2\) and
\(s(Q_1\otimes Q_2)=s_1s_2\), and \(\delta=(d-s)/2\). In particular
\(\delta\) is not multiplicative; the Grothendieck ring of finite nondegenerate
forms under \(\oplus,\otimes\) is \(\Z[\varepsilon]/(\varepsilon^2-1)\) via
\(p+q\varepsilon\), whose only ring homomorphisms to an integral domain are
\(\varepsilon\mapsto\pm1\), i.e.\ \(d\) and \(s\).
\end{corollary}

\begin{proof}
\(s_1s_2=(p_1-q_1)(p_2-q_2)=(p_1p_2+q_1q_2)-(p_1q_2+q_1p_2)\) by Lemma
\ref{lem:whitney}; \(d_1d_2=(p_1+q_1)(p_2+q_2)\) is total dimension. A ring
homomorphism to an integral domain sends \(\varepsilon\) to a root of
\(\varepsilon^2=1\).
\end{proof}

\subsection{The exact characterization of robust finiteness}
The following standard characterization is used throughout. Essential positivity comes in two
strengths, and only the strict one is stable under the compact ideal; the
distinction is the whole content.

\begin{definition}[two essential positivities]\label{def:two-pos}
\(G=G^*\) is \emph{weakly essentially positive} if
\(\specess(G)\subseteq[0,\infty)\), and \emph{strictly essentially positive} if
\(\specess(G)\subseteq[c,\infty)\) for some \(c>0\). Since \(\specess(G)\) is
compact, the latter is equivalent to \(\specess(G)\subseteq(0,\infty)\).
\end{definition}

\begin{theorem}[exact characterization of robust finiteness of the defect]
\label{thm:exact}
For \(G=G^*\in\mathcal B(H)\) the following are equivalent:
\begin{enumerate}
\item[\textup{(i)}] \(G\) is strictly essentially positive;
\item[\textup{(ii)}] \(\delta(G+K)<\infty\) for \emph{every} compact
self-adjoint \(K\);
\item[\textup{(iii)}] there is \(\varepsilon>0\) with
\(\delta(G-\varepsilon I)<\infty\).
\end{enumerate}
\end{theorem}

\begin{proof}
\emph{(i)\(\Rightarrow\)(ii).} By Weyl's theorem on stability of the essential
spectrum under compact perturbations \cite[Thm.\ IV.5.35]{Kato},
\(\specess(G+K)=\specess(G)\subseteq[c,\infty)\). Hence
\(\spec(G+K)\cap(-\infty,0]\) is bounded and discrete (its accumulation points
would lie in \(\specess(G+K)\), disjoint from \((-\infty,0]\)), so finite, of
eigenvalues of finite multiplicity, and \(\delta(G+K)<\infty\) by Lemma
\ref{lem:defect-spectral}.

\emph{(iii)\(\Leftrightarrow\)(i).} \(\delta(G-\varepsilon I)=\dim\ran
E_G(-\infty,\varepsilon)\). If this is finite, \(\spec(G)\cap(-\infty,\varepsilon)\)
is finitely many eigenvalues of finite multiplicity, so
\(\specess(G)\subseteq[\varepsilon,\infty)\): strict. Conversely, if
\(\specess(G)\subseteq[c,\infty)\), take \(\varepsilon=c/2\); the same argument
gives \(\dim\ran E_G(-\infty,c/2)<\infty\).

\emph{(ii)\(\Rightarrow\)(i), by contraposition.} Suppose strictness fails:
there is \(\mu\in\specess(G)\) with \(\mu\leq0\). We construct a compact
self-adjoint \(K\) with \(\delta(G+K)=\infty\).

\emph{Step 1 (orthonormal Weyl sequence with summable error).} For each \(n\),
\(\dim\ran E_G(\mu-8^{-n},\mu+8^{-n})=\infty\) (this is the defining property of
the essential spectrum for self-adjoint operators: if it were finite for some
\(n\), \(\mu\) would be an isolated eigenvalue of finite multiplicity, hence not
in \(\specess\) \cite[\S XIII.4]{ReedSimon}). Choose inductively
\(e_n\in\ran E_G(\mu-8^{-n},\mu+8^{-n})\) of unit norm orthogonal to
\(e_1,\dots,e_{n-1}\) (possible: the range is infinite-dimensional and finitely
many orthogonality constraints). Then \((e_n)\) is orthonormal and
\[
\|(G-\mu)e_n\|^2=\int_{(\mu-8^{-n},\mu+8^{-n})}|\lambda-\mu|^2\,d\|E_\lambda e_n\|^2
\leq 64^{-n},\qquad\text{i.e.}\quad\|(G-\mu)e_n\|\leq 8^{-n}.
\]

\emph{Step 2 (Gram matrix of \(G\)).} For \(n\neq m\), by self-adjointness,
\(|\langle Ge_n,e_m\rangle|=|\langle(G-\mu)e_n,e_m\rangle|\leq 8^{-n}\) and
\(=|\langle e_n,(G-\mu)e_m\rangle|\leq 8^{-m}\), so
\(|\langle Ge_n,e_m\rangle|\leq 8^{-\max(n,m)}\); and
\(\langle Ge_n,e_n\rangle\leq\mu+8^{-n}\leq8^{-n}\) (as \(\mu\leq0\)).

\emph{Step 3 (the perturbation).} Set
\(K:=-\sum_{n\geq1}2^{-n}\langle\cdot,e_n\rangle e_n\): compact (norm limit of
finite-rank), self-adjoint, \(\|K\|=\half\). For
\(x=\sum_{n\geq n_0}\xi_n e_n\) finitely supported in \(\{n\geq n_0\}\),
\[
\langle(G+K)x,x\rangle\leq\sum_n(8^{-n}-2^{-n})|\xi_n|^2
+\sum_{n\neq m}8^{-\max(n,m)}|\xi_n||\xi_m|.
\]
The cross term is bounded by
\[
\sum_{n\neq m}8^{-\max(n,m)}\tfrac{|\xi_n|^2+|\xi_m|^2}{2}
\leq\sum_n|\xi_n|^2 r_n,\qquad
r_n=\sum_{m\neq n}8^{-\max(n,m)}\leq n8^{-n}+\sum_{m>n}8^{-m}\leq(n+1)8^{-n}.
\]
Hence
\(\langle(G+K)x,x\rangle\leq\sum_n\bigl(8^{-n}+(n+1)8^{-n}-2^{-n}\bigr)|\xi_n|^2\).
Since \((n+2)8^{-n}<2^{-n}\) for all \(n\geq4\), taking \(n_0=4\) makes every
coefficient negative: \(G+K\) is negative definite on the
infinite-dimensional \(V_0:=\overline{\operatorname{span}}\{e_n:n\geq4\}\), so
\(\delta(G+K)=\infty\) by Lemma \ref{lem:defect-spectral}.
\end{proof}

\begin{remark}[the design lesson]\label{rem:design}
The only reading of ``essential positivity \(\Rightarrow\) finite defect'' that
survives the compact ideal -- and the essential spectrum lives modulo that ideal
by definition -- is the \emph{strict} one. Any certificate must produce a
\emph{gap} \(c>0\), not merely a sign. The operator \(G=\operatorname{diag}(-1/n)\)
has \(\specess(G)=\{0\}\subseteq[0,\infty)\) (weakly positive) yet
\(\delta=\infty\): by Theorem \ref{thm:exact}, this is the generic shape of
gaplessness. Direction (i)\(\Rightarrow\)(ii) is Weyl's theorem; the converse
is standard (it follows, for instance, from the Weyl--von Neumann theorem); we
included a short proof for completeness.
\end{remark}

\begin{lemma}[gap detector]\label{lem:gap-detector}
If there is an orthonormal sequence \((u_n)\) with
\(\liminf_n\langle Gu_n,u_n\rangle\leq\ell\), then
\(\inf\specess(G)\leq\ell\).
\end{lemma}

\begin{proof}
Suppose \(\specess(G)\subseteq[c,\infty)\) with \(c>\ell\); fix
\(\varepsilon\in(0,c-\ell)\) and \(E:=E_G(-\infty,c-\varepsilon)\), finite-rank
by Theorem \ref{thm:exact}(iii)\(\Leftrightarrow\)(i). With
\(m_0:=\inf\spec(G)>-\infty\),
\[
\langle Gu,u\rangle\geq m_0\|Eu\|^2+(c-\varepsilon)\|(1-E)u\|^2
=(c-\varepsilon)\|u\|^2-(c-\varepsilon-m_0)\|Eu\|^2 .
\]
Since \(u_n\rightharpoonup0\) and \(E\) is finite-rank, \(\|Eu_n\|\to0\), so
\(\liminf\langle Gu_n,u_n\rangle\geq c-\varepsilon>\ell\), a contradiction.
\end{proof}

\section{The Weil--Toeplitz algebra}\label{sec:algebra}

\subsection{The field of windows and the corona}

\begin{definition}[field of windows; \(\W\), \(\Jc\), the symbol]\label{def:field}
A \emph{field of windows} is a sequence of pairs
\((H_n,Q_n)\), \(n\geq2\), where \(H_n\) is a Hilbert space of finite dimension \(d_n\), \(Q_n=Q_n^*\in\mathcal B(H_n)\), and
\(\sup_n\|Q_n\|<\infty\). Set
\[
\W:=\Bigl\{A=(A_n)_n:A_n\in\mathcal B(H_n),\ \|A\|:=\sup_n\|A_n\|<\infty\Bigr\},
\qquad
\Jc:=\{A\in\W:\|A_n\|\to0\}.
\]
With pointwise operations, \(\W\) is a \(C^*\)-algebra (the \(\ell^\infty\)
product of the field) and \(\Jc\) is a closed two-sided ideal (the
\emph{compact ideal of the field}: it is the closure of the finitely supported
sections). The \emph{corona} is \(\Cor:=\W/\Jc\) and the \emph{symbol} of a
section is its class \(\sigma(A):=A+\Jc\in\Cor\).
\end{definition}

\begin{remark}
In the block-diagonal operator \(\bigoplus_n A_n\) on \(\bigoplus_n H_n\), a
section lies in \(\Jc\) iff the operator is compact; \(\Cor\) is the image of
the section algebra in the Calkin algebra of \(\bigoplus_n H_n\). The corona is
literally ``the family modulo compacts''. The name Weil--Toeplitz reflects the
analogy with the Toeplitz symbol: the essential spectrum of \(T_\phi\) is the
range of the symbol -- the ``limits of the windows'' of the operator
\cite{BottcherSilbermann,RabinovichRochSilbermann}.
\end{remark}

\begin{lemma}[norm of the symbol]\label{lem:symbol-norm}
\(\|\sigma(A)\|_{\Cor}=\limsup_n\|A_n\|\).
\end{lemma}

\begin{proof}
(\(\leq\)) For fixed \(N\) the truncation
\(A^{(N)}:=(A_2,\dots,A_N,0,0,\dots)\in\Jc\), so
\(\|\sigma(A)\|\leq\|A-A^{(N)}\|=\sup_{n>N}\|A_n\|\to\limsup_n\|A_n\|\). (\(\geq\))
For every \(J\in\Jc\),
\(\|A+J\|=\sup_n\|A_n+J_n\|\geq\limsup_n\|A_n+J_n\|=\limsup_n\|A_n\|\) since
\(\|J_n\|\to0\); take the infimum over \(J\).
\end{proof}

\begin{lemma}[symbol spectrum is the limit set of window spectra]
\label{lem:symbol-spectrum}
For \(A=(A_n)\in\W\) with each \(A_n\) self-adjoint,
\[
\spec_{\Cor}(\sigma(A))=\Lambda(A):=\bigl\{\lambda\in\R:\exists\,n_k\to\infty,\
\lambda_k\in\spec(A_{n_k}),\ \lambda_k\to\lambda\bigr\}.
\]
\end{lemma}

\begin{proof}
\(\sigma(A)\) is self-adjoint, so its spectrum is real; \(\Lambda(A)\) is closed
and bounded.

(\(\subseteq\) via contrapositive) If \(\sigma(A)-\lambda\) is invertible in
\(\Cor\), there is \(B\in\W\) with \((A-\lambda)B=1+J\), \(B(A-\lambda)=1+J'\),
\(J,J'\in\Jc\). For large \(n\), \(\|J_n\|,\|J'_n\|\leq\half\), so
\((A_n-\lambda)B_n\) and \(B_n(A_n-\lambda)\) are invertible, whence
\(A_n-\lambda\) is invertible with \(\|(A_n-\lambda)^{-1}\|\leq2\|B\|\). For
self-adjoint \(A_n\),
\(\|(A_n-\lambda)^{-1}\|=\dist(\lambda,\spec A_n)^{-1}\), so
\(\dist(\lambda,\spec A_n)\geq(2\|B\|)^{-1}\) for all large \(n\): no
\(\lambda_k\in\spec(A_{n_k})\) can converge to \(\lambda\), i.e.\
\(\lambda\notin\Lambda(A)\).

(\(\supseteq\) via contrapositive) If \(\lambda\notin\Lambda(A)\), there are
\(\varepsilon>0\) and \(N\) with \(\dist(\lambda,\spec A_n)\geq\varepsilon\) for
all \(n\geq N\) (else a violating subsequence would, by boundedness of the
spectra, produce a point of \(\Lambda\) within \(\varepsilon\) of \(\lambda\)).
Put \(B_n:=(A_n-\lambda)^{-1}\) for \(n\geq N\) (\(\|B_n\|\leq\varepsilon^{-1}\))
and \(B_n:=0\) for \(n<N\); then \(B\in\W\) and \((A-\lambda)B-1\),
\(B(A-\lambda)-1\) are finitely supported, hence in \(\Jc\), so \(\sigma(B)\)
inverts \(\sigma(A)-\lambda\).
\end{proof}

\begin{corollary}[positivity of the symbol is the uniform quantifier]
\label{cor:master-quantifier}
For \(A=A^*\):
\begin{itemize}
\item \(\spec(\sigma(A))\subseteq[0,\infty)\) \emph{(weak symbol positivity)}
\(\iff\) \(\min\spec(A_n)\geq-\varepsilon_n\) with \(\varepsilon_n\to0\);
\item \(\spec(\sigma(A))\subseteq[c,\infty)\), \(c>0\) \emph{(strict symbol
positivity)} \(\iff\) for all \(\varepsilon>0\) there is \(N\) with
\(\spec(A_n)\subseteq(c-\varepsilon,\infty)\) for all \(n\geq N\) -- a gap
\emph{uniform over windows, eventually}.
\end{itemize}
\end{corollary}

\begin{proof}
Read Lemma \ref{lem:symbol-spectrum} twice.
\end{proof}

This is the structural observation that orients the paper: the master
quantifier (mean over windows versus uniformly over windows) becomes membership
in \(c_0\) versus \(\ell^\infty\). Corollary \ref{cor:master-quantifier} is the
statement ``uniform in every window'' written as positivity of an element of a
\(C^*\)-algebra.

\subsection{The field of Weil}
We now describe the field carrying the Weil form. The structural inputs
are isolated as hypotheses.

\begin{definition}[the Weil field at a resolution calendar]\label{def:weil-field}
Fix \(\Delta\in[\tfrac12,1]\) (so that every ordinate \(\gamma\geq\tfrac32\) lies
in at least one and at most three windows) and a nondecreasing \emph{resolution calendar}
\(a\colon[2,\infty)\to[1,\infty)\) (the log-support budget of the tests per
window; the natural choice is \(a(T)\asymp\log T\), the resolution scale of
individual zeros). For each \(n\): \(H_n\) is a finite-dimensional subspace of
the test class, concentrated in the window \(|\operatorname{Im}s-n|\leq\Delta\)
with multiplicative support \(\subseteq[e^{-a(n)},e^{a(n)}]\), endowed with the
\emph{online-sea normalization} (the inner product is the spectral Gram form of
the online zeros of the window, normalized to mass \(1\) per zero); and \(Q_n\)
is the compression of the Weil form to \(H_n\) in that normalization. The
\emph{autonomous} section \(\Qaut=(\Qaut_n)\) is the polar/archimedean-plus-near-
prime prediction of each window; that it is computable
without zero positions is part of Hypothesis \ref{hyp:autonomy}.
\end{definition}

The following are hypotheses, not proved here or elsewhere in the series.

\begin{assumption}[block structure]\label{hyp:blocks}
Each \(Q_n\) decomposes, up to controlled couplings, as an online-sea block
(\(\approx\) identity after normalization) plus one hyperbolic plane
\(\left(\begin{smallmatrix}0&\theta\\\theta&0\end{smallmatrix}\right)\) for each
visible off-critical quadruple in the window. \emph{(For the unwindowed form
on even test functions, one Lorentzian block of signature \((1,1)\) per
off-critical quartet is \cite[Theorem~3.1]{P18}; the windowed statement is not
proved there.)}
\end{assumption}

\begin{assumption}[autonomy of the value]\label{hyp:autonomy}
For tests \(f,g\) of the class, \(Q_n(f,g)\) is computable from the arithmetic
blocks of the explicit formula (primes, pole, gamma factor) without any zero
position. \emph{(For the unnormalized form this is Weil's explicit formula
\cite{Weil1952}; the hypothesis is that it persists in the online-sea
normalization.)}
\end{assumption}

\begin{assumption}[off-critical count]\label{hyp:spectral}
Let \(m\in\N\cup\{\infty\}\) be the number of off-critical quadruples; thus
\(\mathrm{RH}\iff m=0\), and each quadruple contributes one hyperbolic plane.
\emph{(The first statement is the definition of \(m\). On the even test space,
each off-critical quartet contributes one Lorentzian block and the negative index
is at most \(m\) \cite[Theorem~3.1, Corollary~6.5]{P18}. On the full test space,
if the off-line zeros form \(m<\infty\) quartets, the negative index is \(2m\)
(one plane per pair \(\{\rho,1-\bar\rho\}\)) and equals the number of
eigenvalues of the scaling generator in the upper half-plane
\cite[Proposition~5.3]{P25}; a window at positive height contains one such pair
per quadruple. We do not use this.)}
\end{assumption}

\begin{remark}[status of the field]\label{rem:field-status}
That the compressions exist with the stated normalization, and that the section
is bounded, is not proved here. The uniform
boundedness \(\sup_n\|Q_n\|<\infty\) in the online-sea normalization requires
control of the per-window pair-correlation corrections; we treat this as an
axiom \textbf{(S1)} of the model \(\mathfrak W\) below
(per-window pair-correlation boundedness, not proved here), inside which all
theorems of Section \ref{sec:trichotomy} are proved. This separates theorems in
the axiomatized model from their transfer to the genuine Weil instance (Section
\ref{sec:honest}).
\end{remark}

\section{The spectral model and the trichotomy}\label{sec:trichotomy}

\subsection{The model}
We axiomatize the spectral structure of the windowed Weil form, distilled from
Hypotheses \ref{hyp:blocks}--\ref{hyp:spectral}.

\begin{definition}[configuration and visibility]\label{def:config}
A \emph{configuration} is an at most countable set
\(Z=\{(\delta_j,\gamma_j)\}_{j\in J}\) with \(\delta_j\in(0,\half)\), \(\gamma_j\geq\tfrac32\),
\(\gamma_j\to\infty\) if \(|J|=\infty\), and finitely many \(j\) per window;
\(m:=|J|\). (\(Z=\emptyset\) is ``RH''.) The \emph{visibility} of quadruple
\(j\) at calendar \(a\) is \(\theta_j:=\theta(\delta_j,\gamma_j)\in[0,1]\),
where \(\theta\) is a fixed function with
\begin{itemize}
\item[\textup{(V1)}] \(\theta\) continuous, nondecreasing in \(\delta\), and
\(\theta(\delta,\gamma)\to0\) as \(\delta\,a(\gamma)\to0\);
\item[\textup{(V2)}] there is \(c_0>0\) with \(\theta(\delta,\gamma)\geq c_0\)
when \(\delta\,a(\gamma)\geq1\);
\item[\textup{(V3)}] \emph{(optional; assumed only where stated)} there is
\(t_0\in(0,1]\) with \(\theta(\delta,\gamma)=0\) whenever
\(\delta\,a(\gamma)<t_0\).
\end{itemize}
Quadruples with \(\delta_j\,a(\gamma_j)\geq1\) are \emph{supra-resolution (fat)};
with \(\delta_j\,a(\gamma_j)<1\), \emph{sub-resolution (thin)}; a thin quadruple
is \emph{visible} if \(\theta_j>0\) and \emph{blind} if \(\theta_j=0\).
\end{definition}

\begin{remark}[origin of \(\theta\)]\label{rem:theta}
(V1)--(V2) are motivated by a Paley--Wiener estimate: a
test with multiplicative support in \([e^{-a},e^a]\) has
\(\hat f(s)=\int_{-a}^{a}F(u)e^{(s-1/2)u}\,du\), so over a band of width
\(2\delta\),
\[
\hat f(\tfrac12+\delta+i\gamma)\,\overline{\hat f(\tfrac12-\delta+i\gamma)}
=\bigl|\hat f(\tfrac12+i\gamma)\bigr|^2
+O\bigl(\delta\,a\,e^{a/2}\|F\|_{L^1}\sup_{\text{band}}|\hat f|\bigr).
\]
Thus a quadruple with \(\delta a\to0\) differs from an online pair by
\(O(\delta a)\to0\) per unit mass (V1); for \(\delta a\geq1\) an aligned test
realizes a sign opposition of size \(\geq c_0\) (V2), if the class is closed
under multiplication by Mellin polynomials. This is a heuristic, not a proof;
the exact behaviour in the transition \(\delta a\asymp1\) is not used: no
theorem below uses more than (V1)--(V3).
\end{remark}

\begin{definition}[the model \(\mathfrak W(Z)\)]\label{def:model}
Given a configuration \(Z\) and calendar \(a\), set
\[
H_n:=\C^{d_n}\oplus\bigoplus_{j\in J_n^{\mathrm{vis}}}\C^2,\qquad
Q_n:=(I_{d_n}+R_n)\oplus\bigoplus_{j\in J_n^{\mathrm{vis}}}
\begin{pmatrix}0&\theta_j\\\theta_j&0\end{pmatrix}+C_n,
\]
where \(d_n\asymp\Delta\log n\),
\(J_n^{\mathrm{vis}}:=\{j:|\gamma_j-n|\leq\Delta,\ \theta_j>0\}\), and
\begin{itemize}
\item[\textup{(S1)}] \(\|R_n\|\leq r<\half\) for all \(n\) \emph{(sea
corrections; not proved for the Weil instance, Remark \ref{rem:field-status})};
\item[\textup{(S2)}] the hyperbolic blocks carry weight \(\theta_j\) with
(V1)--(V2) \emph{(Remark \ref{rem:theta} for the Weil instance)};
\item[\textup{(S3)}] \(\|C_n\|\leq\omega_n\) with \(\omega_n\to0\)
\emph{(cross-couplings between separated windows; not proved for the Weil
instance)}.
\end{itemize}
The autonomous section is \(\Qaut_n:=(I_{d_n}+R_n)\oplus
I_{2|J_n^{\mathrm{vis}}|}\). Sub-resolution quadruples with \(\theta_j=0\) do not
appear in \(H_n\); their entire effect is absorbed into \(R_n/C_n\).
\end{definition}

All theorems in this section are \emph{unconditional inside} \(\mathfrak W\);
their transport to \(\zeta\) carries the statuses of (S1)--(S3) and is discussed
in Section \ref{sec:honest}.

\subsection{The stratification theorems}
Throughout, \(\sigma(Q)\in\Cor\), \(\Lambda(Q)=\spec(\sigma(Q))\), and
\(\vartheta_n:=\max\{\theta_j:j\in J_n^{\mathrm{vis}}\}\) (\(:=0\) if empty).

\begin{lemma}[per-window spectrum]\label{lem:window-spec}
For all \(n\),
\(\spec(Q_n)\subseteq\bigl([1-r,1+r]\cup\{\pm\theta_j\}_{j\in J_n^{\mathrm{vis}}}
\bigr)+[-\omega_n,\omega_n]\). In particular
\(\lambda_{\min}(Q_n)\geq-\vartheta_n-\omega_n\); and if
\(J_n^{\mathrm{vis}}\neq\emptyset\),
\(\lambda_{\min}(Q_n)\leq-\vartheta_n+\omega_n\).
\end{lemma}

\begin{proof}
Without \(C_n\), \(Q_n\) is a direct sum: the sea block has spectrum in
\([1-r,1+r]\), and each
\(\left(\begin{smallmatrix}0&\theta\\\theta&0\end{smallmatrix}\right)\) has
eigenvalues \(\pm\theta\). The self-adjoint perturbation \(C_n\) moves each
eigenvalue by at most \(\|C_n\|\leq\omega_n\) (finite-dimensional Weyl
inequality \cite{Bhatia}). The upper bound: the vector \(v=(1,-1)/\sqrt2\) of
the \(\vartheta_n\)-block gives
\(\lambda_{\min}(Q_n)\leq\langle Q_nv,v\rangle\leq-\vartheta_n+\omega_n\).
\end{proof}

\begin{theorem}[fat \(\Rightarrow\) the symbol sees them]\label{thm:fat}
If there are \(c>0\) and a subsequence \(j_k\) with \(\theta_{j_k}\geq c\) and
\(\gamma_{j_k}\to\infty\), then
\(\spec(\sigma(Q))\cap(-\infty,-c/2]\neq\emptyset\). In particular the symbol is
not even weakly positive: the negativity of supra-resolution quadruples
\emph{survives the quotient by the compact ideal}.
\end{theorem}

\begin{proof}
Let \(n_k\) be the window of \(\gamma_{j_k}\) (\(n_k\to\infty\)). By Lemma
\ref{lem:window-spec}, \(\lambda_{\min}(Q_{n_k})\leq-c+\omega_{n_k}\). Since
\(\omega_n\to0\) and the spectra are uniformly bounded, a subsequence of
\(\lambda_{\min}(Q_{n_k})\) converges to a point \(\leq-c\), which lies in
\(\Lambda(Q)\) by Lemma \ref{lem:symbol-spectrum}.
\end{proof}

\begin{theorem}[thin visible \(\Rightarrow\) negatives live exactly at the
boundary \(0\)]\label{thm:thin}
Suppose \(m=\infty\), \(J_n^{\mathrm{vis}}\neq\emptyset\) for infinitely many
\(n\), and \(\vartheta_n\to0\). Then:
\begin{enumerate}
\item[\textup{(i)}] \(0\in\spec(\sigma(Q))\); if moreover \(\vartheta_n>\omega_n\)
for infinitely many populated windows, it is attained as a limit of genuinely
\emph{negative} \(\lambda_{\min}(Q_n)\) over those windows;
\item[\textup{(ii)}] \(\spec(\sigma(Q))\subseteq\{0\}\cup[1-r,1+r]\): weak symbol
positivity \emph{holds};
\item[\textup{(iii)}] strict symbol positivity \emph{fails}.
\end{enumerate}
The negatives accumulate exactly at the single point \(0\) -- the point
separating weak from strict positivity, the very point of Theorem
\ref{thm:exact}.
\end{theorem}

\begin{proof}
(i) Over the infinitely many populated windows,
\(\lambda_{\min}(Q_n)\in[-\vartheta_n-\omega_n,-\vartheta_n+\omega_n]\to0\)
(Lemma \ref{lem:window-spec}, \(\vartheta_n,\omega_n\to0\)); when
\(\vartheta_n>\omega_n\) these are genuinely negative, and the limit \(0\) lies
in \(\Lambda(Q)\) (Lemma \ref{lem:symbol-spectrum}). (ii) Every eigenvalue of
\(Q_n\) lies in
\([-\vartheta_n-\omega_n,\vartheta_n+\omega_n]\cup[1-r-\omega_n,1+r+\omega_n]\);
its limit points lie in \(\{0\}\cup[1-r,1+r]\). (iii) Immediate from (i) and
Corollary \ref{cor:master-quantifier}.
\end{proof}

\begin{theorem}[strict symbol positivity \(\iff\) finite \emph{visible} count]
\label{thm:strict}
Strict symbol positivity at calendar \(a\)
holds iff \(\#\{n:J_n^{\mathrm{vis}}\neq\emptyset\}<\infty\), i.e.\
\(m_a^{\mathrm{vis}}:=\#\{j:\theta_j>0\}<\infty\). In particular strict positivity says \emph{nothing} about
\(m-m_a^{\mathrm{vis}}\), the sub-resolution count.
\end{theorem}

\begin{proof}
(\(\Leftarrow\)) If \(J_n^{\mathrm{vis}}=\emptyset\) for \(n\geq N\),
then \(\spec(Q_n)\subseteq[1-r-\omega_n,1+r+\omega_n]\) for \(n\geq N\), so
\(\Lambda(Q)\subseteq[1-r,1+r]\) with gap \(c=1-r>0\) (Corollary
\ref{cor:master-quantifier}). (\(\Rightarrow\)) If infinitely many windows are
populated, each carries an eigenvalue
\(\mu_n\in[-\vartheta_n-\omega_n,-\vartheta_n+\omega_n]\) (Lemma
\ref{lem:window-spec}). The sequence \((\mu_n)\) over the populated windows is
bounded; any accumulation point \(\ell\) satisfies
\(\ell\leq\limsup_n(-\vartheta_n+\omega_n)\leq\limsup_n\omega_n=0\), and
\(\ell\in\Lambda(Q)\) by Lemma \ref{lem:symbol-spectrum}: strict positivity
fails, \emph{regardless of the comparison between \(\vartheta_n\) and
\(\omega_n\)}.
\end{proof}

\begin{remark}[sub-noise windows]\label{rem:strict-corrected}
Theorem \ref{thm:strict} admits no exception for visible windows with
\(\vartheta_n\leq\omega_n\). For example, with \(\omega_n=1/n\) and one visible
quadruple in each even window with \(\vartheta_n=1/(2n)\), the perturbed block
eigenvalues lie in \([-3/(2n),1/(2n)]\to0\), so \(0\in\spec(\sigma(Q))\) and
there is no gap. If one wishes to discard visible quadruples below the noise
level, this must be done in the definition of \(J_n^{\mathrm{vis}}\) (a
threshold \(\theta_j>\omega\)), not in the theorem.
\end{remark}

\begin{corollary}[the gap controls the visible count quantitatively]
\label{cor:gap-quant}
Let \(c>0\) and assume \(\omega_n<c\) for \(n\geq N_0\). If
\(\spec(Q_n)\subseteq[c,\infty)\) for \(n\geq N_0\), then
\(J_n^{\mathrm{vis}}=\emptyset\) for \(n\geq N_0\), and
\(m_a^{\mathrm{vis}}\leq\sum_{n<N_0}|J_n^{\mathrm{vis}}|<\infty\). If moreover
\(Z\) consists of quadruples of zeros of \(\zeta\), the sum is
\(\ll N_0\log N_0\) by \cite[Thm.~5.8]{IwaniecKowalski}.
\end{corollary}

\begin{proof}
If \(J_n^{\mathrm{vis}}\neq\emptyset\) for some \(n\geq N_0\), Lemma
\ref{lem:window-spec} gives \(\lambda_{\min}(Q_n)\leq-\vartheta_n+\omega_n
\leq\omega_n<c\), a contradiction. Finiteness of the sum holds since each window carries finitely
many quadruples (Definition \ref{def:config}). For zeros of \(\zeta\), each window \(|\gamma-n|\leq\Delta\) contains
\(\ll\log(2+n)\) zeros by the zero-counting estimate \(N(t+1)-N(t)\ll\log(2+t)\)
\cite[Thm.~5.8]{IwaniecKowalski}; sum over \(n<N_0\).
\end{proof}

\begin{theorem}[sub-resolution \(\Rightarrow\) the symbol is blind]
\label{thm:blind}
Let \(Z\) have \(m=\infty\) and \(J_n^{\mathrm{vis}}=\emptyset\) for all \(n\)
(all quadruples sub-resolution with \(\theta_j=0\)). Then, as sections over the
same field,
\[
Q-\Qaut\in\Jc,\qquad\text{hence}\qquad\sigma(Q)=\sigma(\Qaut)\ \text{in }\Cor,
\]
and \(\sigma(Q)\) is strictly positive, with
\(\spec\sigma(Q)\subseteq[1-r,1+r]\). Strict symbol positivity does \emph{not}
imply \(m<\infty\).
\end{theorem}

\begin{proof}
Since \(J_n^{\mathrm{vis}}=\emptyset\), Definition \ref{def:model} gives
\(Q_n-\Qaut_n=C_n\), so \(\|Q_n-\Qaut_n\|\leq\omega_n\to0\) by (S3); hence
\(Q-\Qaut\in\Jc\). Then \(\Qaut_n=I_{d_n}+R_n\) has spectrum in \([1-r,1+r]\)
by (S1), and Lemma \ref{lem:symbol-spectrum} gives
\(\spec\sigma(Q)=\spec\sigma(\Qaut)\subseteq[1-r,1+r]\), a gap \(1-r>0\),
while \(m=\infty\).
\end{proof}

If, outside the model, each quadruple contributes one hyperbolic plane to the
unrestricted form (with tests of support exceeding the calendar), Lemma
\ref{lem:mono-add}(ii) gives defect \(\geq m=\infty\) for that form; this
observation is not part of Theorem \ref{thm:blind}.

\begin{theorem}[calendar no-go]\label{thm:nogo}
The finiteness \(m<\infty\) is a boundary event over the resolution filtration:
\begin{enumerate}
\item[\textup{(a)}] Assume \textup{(V3)}. For every calendar \(a\), the
configuration given by \(\gamma_j:=j\), \(j\geq2\), and
\[
\delta_j:=\min\bigl(\tfrac14,\,t_0/(2j\,a(j))\bigr)
\]
has \(m=\infty\) and
\(\theta_j=0\) for all \(j\); hence \(J_n^{\mathrm{vis}}=\emptyset\) and, by
Theorem \ref{thm:blind}, \(\sigma(Q)=\sigma(\Qaut)\) is strictly positive.
Without \textup{(V3)}, every configuration with \(\theta_j>0\) for infinitely
many \(j\) and \(\vartheta_n\to0\) is thin-visible, and by Theorem
\ref{thm:thin} its symbol is weakly but not strictly positive.
\item[\textup{(b)}] For every configuration and quadruple, some calendar sees it
(\(a(\gamma_j)\geq1/\delta_j\)). Hence with
\(m_a:=\#\{j:\delta_j a(\gamma_j)\geq1\}\), \(\;m=\sup_a m_a\).
\item[\textup{(c)}] Consequently \(m<\infty\iff\exists M\,\forall a:m_a\leq M\),
a \(\Sigma_2\) statement over the filtration of resolutions; under
\textup{(V3)}, strict positivity of the symbol at a fixed calendar does not
certify \(m<\infty\).
\end{enumerate}
\end{theorem}

\begin{proof}
(a) \(\delta_j a(\gamma_j)\leq t_0/(2j)<t_0\), so \(\theta_j=0\) by (V3); the
configuration is infinite with at most three quadruples per window; apply Theorem
\ref{thm:blind}. The statement without (V3) is Theorem \ref{thm:thin}. (b)
Monotonicity of \(a\mapsto m_a\) and the pointwise observation. (c) Rewriting of
(b) together with (a).
\end{proof}

\begin{remark}[the trichotomy]\label{rem:trichotomy}
The dichotomy ``symbol or ideal'' is false; it is a trichotomy stratified by
\(\delta_j a(\gamma_j)\):
\[
\begin{array}{lll}
\text{fat }(\delta_j a(\gamma_j)\geq1\text{ i.o.}) &
\text{in the symbol: mass }\leq-c/2 & \text{Thm.\ \ref{thm:fat}}\\
\text{thin visible }(\theta_j\to0^+) &
\text{at the boundary }0\in\spec\sigma(Q) & \text{Thm.\ \ref{thm:thin}}\\
\text{zero visibility }(\theta_j=0\ \forall j;\text{ under (V3), }\delta_ja(\gamma_j)<t_0) &
\text{in the ideal: }\sigma(Q)=\sigma(\Qaut) & \text{Thm.\ \ref{thm:blind}}
\end{array}
\]
The strata are relative to the calendar (Theorem \ref{thm:nogo}); within
\(\mathfrak W\), it is the \emph{intermediate (thin-visible) stratum} that lives
on the symbol/compact boundary, accumulating at \(0\) -- the unique point
separating weak from strict positivity, which by Theorem \ref{thm:exact} is the
only point on which robust finiteness of the defect depends -- while the fat
stratum sits in the negative interior of the symbol and the sub-resolution
stratum in the ideal. The slogan ``the zeros live on the boundary'' is therefore
a theorem-in-model for the intermediate stratum only, not a global statement
about the zeros of \(\zeta\). The master quantifier acquires an
exact \(C^*\)-algebraic direction: it is the point \(0\in\spec_{\Cor}\).
\end{remark}

\subsection{Synergy of the two axes}
Symbol positivity and the replication axiom (R) combine cleanly.

\begin{theorem}[replication \(\Rightarrow\) the symbol sees]\label{thm:synergy}
Suppose \(Z\) satisfies the quantitative replication axiom \emph{(R)}: if
\(m\geq1\), there is a quadruple \((\delta_0,\gamma_0)\in Z\) and heights
\(\tau_k\to\infty\) such that \(Z\) contains quadruples
\((\delta_k',\gamma_0+\tau_k)\) with \(\delta_k'\geq\delta_0/2\). Let \(a\) be any
calendar with \(a(\gamma)\to\infty\). Then \(m\geq1\) forces infinitely many fat
quadruples, hence
\(\spec(\sigma(Q))\cap(-\infty,-c_0/2]\neq\emptyset\). Contrapositively,
\[
\text{weak symbol positivity }\wedge\text{ (R)}\ \Longrightarrow\ m=0
\text{ in }\mathfrak W.
\]
For \(\zeta\) this gives RH only if the windows of \(\zeta\) satisfy the model
axioms (S1)--(S3), which are not proved.
\end{theorem}

\begin{proof}
The replicas have
\(\delta_k'a(\gamma_0+\tau_k)\geq(\delta_0/2)a(\gamma_0+\tau_k)\to\infty\), hence
are eventually supra-resolution with visibility \(\geq c_0\) by (V2). Theorem
\ref{thm:fat} applies. The contrapositive uses only weak positivity since
Theorem \ref{thm:fat} violates even weak positivity.
\end{proof}

The thin stratum, which defeats symbol positivity alone (Theorems
\ref{thm:blind}, \ref{thm:nogo}), is a stratum that replication does
not produce: replicas inherit a fixed \(\delta_0\), and every unbounded
calendar eventually sees them. We record the missing piece this leaves open.

\begin{conjecture}[resolution repulsion]
\label{conj:lp134}
There exist a calendar \(a\) (e.g.\ \(a(\gamma)=C\log\gamma\)) and \(c>0\) such
that every nontrivial zero \(\rho=\half+\delta+i\gamma\) of \(\zeta\) with
\(\delta\neq0\) satisfies \(|\delta|\geq c/a(\gamma)\). \emph{(``Zeros cannot be
asymptotically sub-resolution.'')}
\end{conjecture}

Under Conjecture \ref{conj:lp134}, with the calendar \(a/c\) in place of \(a\)
(we may take \(c\leq1\)), the thin and sub-resolution strata are empty, and by
Theorems \ref{thm:fat}, \ref{thm:strict} symbol positivity becomes a single
\(C^*\)-statement: \(m<\infty\iff\) strict symbol positivity at one calendar.
Conjecture \ref{conj:lp134} is weaker than RH (a quadruple at
\(\delta=\tfrac14\) satisfies it and violates RH). Classical zero-free regions
concern the neighbourhood of \(\operatorname{Re}s=1\), and density theorems
bound how many zeros lie off the line, not how close they are to it.

\section{The localized signature and the logistic law}\label{sec:logistic}

We now develop the second, independent contribution: an exact logistic law for
the tensor tower of indices, and its application to a localized quantity built
from a doubled form \(Q_2\) defined on the zero side. We work in
non-standard coordinates, using a non-principal ultrafilter
\(\mathcal U\) on \(\N\), the ultrapower \({}^*\R=\R^\N/\mathcal U\), and the
standard part \(\sst\) on the limited elements \cite{Robinson}.

\subsection{The logistic law}

\begin{theorem}[exact logistic law of the tensor tower]\label{thm:logistic}
Let \(\bar Q\) be a nondegenerate hermitian form of signature \((p,q)\) on a
space of dimension \(\bar n=p+q\), and set
\(x_k:=\negind(\bar Q^{\otimes 2^k})/\bar n^{2^k}\). Then for all \(k\geq0\), as
an identity of rationals,
\[
x_{k+1}=2x_k(1-x_k),\qquad 1-2x_k=(1-2x_0)^{2^k}.
\]
Consequently, for any scheme of forms with \(\sigma^0:=\sst([x_0(N)])\),
\[
\sigma^{(k)}:=\sst([x_k(N)])=\tfrac12\bigl(1-(1-2\sigma^0)^{2^k}\bigr),
\qquad
\lim_{k\to\infty}\sigma^{(k)}=
\begin{cases}0,&\sigma^0\in\{0,1\},\\[2pt]\tfrac12,&\sigma^0\in(0,1).\end{cases}
\]
\end{theorem}

\begin{proof}
By Lemma \ref{lem:whitney} applied to
\(\bar Q^{\otimes 2^{k+1}}=(\bar Q^{\otimes 2^k})\otimes(\bar Q^{\otimes 2^k})\),
with signature \((P_k,K_k)\), \(P_k+K_k=\bar n^{2^k}\),
\[
K_{k+1}=2P_kK_k,\qquad P_{k+1}=P_k^2+K_k^2.
\]
Dividing by \(\bar n^{2^{k+1}}=(\bar n^{2^k})^2\) gives
\(x_{k+1}=2(P_k/\bar n^{2^k})(K_k/\bar n^{2^k})=2x_k(1-x_k)\). Then
\(1-2x_{k+1}=1-4x_k+4x_k^2=(1-2x_k)^2\), so \(1-2x_k=(1-2x_0)^{2^k}\). Each
\(x_k(N)\in[0,1]\) is limited; \(\sst\) is a ring homomorphism on limited
elements, so the polynomial identity \(1-2x_{k+1}=(1-2x_k)^2\) passes to
standard parts for each fixed \(k\) (\(k\) standard; the identity is used
finitely often). Limits: if \(\sigma^0\in(0,1)\) then \(|1-2\sigma^0|<1\) and
\((1-2\sigma^0)^{2^k}\to0\); if \(\sigma^0=0\), \(\sigma^{(k)}\equiv0\); if
\(\sigma^0=1\), \(1-2\sigma^0=-1\) and \((-1)^{2^k}=1\) for \(k\geq1\), so
\(\sigma^{(k)}=0\) -- consistent, since the tensor square of a negative-definite
form is positive-definite.
\end{proof}

\begin{remark}[the global signature collapses]\label{rem:collapse}
If the index of a window is normalized by the full window dimension, which tends
to infinity, a defect of bounded cardinality has density \(0\); this motivates
the localized construction below, where the cell dimension is bounded. We do not
use this remark.
\end{remark}

\subsection{The localized signature}

\emph{The doubled form.} Count nontrivial zeros with multiplicity, and for a pair
\(P=(\rho_1,\rho_2)\) of nontrivial zeros put \(\sigma=(\rho_1+\rho_2)/2\),
\(\tau=(\rho_1-\rho_2)/2\) and \(\iota(P):=(1-\bar\rho_1,1-\bar\rho_2)\). By
the functional equation and the reflection principle, \(\iota\) is an
involution of the set of pairs preserving multiplicities; it preserves the
ordinates, \(\operatorname{Im}\sigma(\iota P)=\operatorname{Im}\sigma(P)\),
\(\operatorname{Im}\tau(\iota P)=\operatorname{Im}\tau(P)\), and
\(0<\operatorname{Re}\sigma<1\), \(|\operatorname{Re}\tau|<\half\). Let
\(\mathcal E\) be the space of entire functions \(F\) with
\(\sup_{|\operatorname{Re}z|\leq1}(1+|\operatorname{Im}z|)^3|F(z)|<\infty\).
For \(\hat g,\hat\varphi\in\mathcal E\) let \(h_{g,\varphi}\) denote the
function \(\hat h_{g,\varphi}(\rho_1,\rho_2):=\half\hat g(\sigma)\hat\varphi(\tau)\)
on pairs of zeros, let \(\mathcal H\) be the span of these functions, and set,
for \(h,h'\in\mathcal H\),
\begin{equation}\label{eq:speca}
Q_2(h,h')=\sum_{(\rho_1,\rho_2)}\hat h(\rho_1,\rho_2)\,
\overline{\hat h'(\iota(\rho_1,\rho_2))} .
\end{equation}
The series converges absolutely: by the elementary inequality
\eqref{eq:elem} in the proof of Lemma \ref{lem:tail} (with \(\eta=\eta'=0\)),
\(|\hat h_{g,\varphi}(P)\hat h_{g',\varphi'}(\iota P)|\leq
C(1+|\gamma_1|)^{-3}(1+|\gamma_2|)^{-3}\), and
\(\sum_\rho(1+|\gamma|)^{-3}<\infty\) by the Riemann--von Mangoldt formula.
Reindexing \(P\mapsto\iota P\) shows \(Q_2(h',h)=\overline{Q_2(h,h')}\), so
\(Q_2\) is hermitian. Pairs fixed by \(\iota\) (both members on the critical
line) contribute \(|\hat h(P)|^2\geq0\); a free orbit \(\{P,\iota P\}\)
contributes \(2\operatorname{Re}[\hat h(P)\overline{\hat h(\iota P)}]\), a
hyperbolic plane of signature \((1,1)\) when \(\hat h(P)\) and
\(\hat h(\iota P)\) are independent functionals. Under RH every pair is fixed by
\(\iota\), and \(Q_2(h,h)=\sum_P|\hat h(P)|^2\geq0\). The form \(Q_2\) is
defined here on the zero side; we do not establish an expression of \(Q_2\)
through primes (a two-variable analogue of Weil's explicit formula).

\begin{definition}[micro-resolution cell]\label{def:cell}
Fix a \emph{profile} \(p=(D,w,M,b,\eta,\eta')\) with \(D\in\N\) (degree), \(w>0\)
(width), \(M\in\N\) with \(M\geq D+3\) (decay order), \(b\in\mathbb Q_{>0}\)
(caliber), \((\eta,\eta')\in\mathbb Q^2\) (centres). The \emph{cell} is
\[
V_p:=\operatorname{span}\bigl\{h_{g_a,\varphi_e}:0\leq a,e\leq D\bigr\},\quad
\hat g_a(\sigma):=B^{(M)}_w(\sigma-i\eta)\sigma^a,\quad
\hat\varphi_e(\tau):=B^{(M)}_w(\tau-i\eta')\tau^e,
\]
where \(B^{(M)}_w\not\equiv0\) is the Mellin transform of a fixed even bump
\(g_0^{(M,w)}\in C_c^\infty(\R_{>0})\) with
\[
|B^{(M)}_w(z)|\leq C_M e^{a_0|\operatorname{Re}z-1/2|}
(1+|\operatorname{Im}z|/w)^{-M}
\tag{$\ast$}
\]
(Paley--Wiener; \(a_0=a_0(M,w)\) the fixed log-support). Then
\(\hat g_a,\hat\varphi_e\in\mathcal E\), \(\dim V_p\leq(D+1)^2\), and \(V_p\) is
test-accessible: the profile mentions no zeros, and each \(\hat g_a\) is the
Mellin transform of a function in \(C_c^\infty(\R_{>0})\) (multiplication of a
Mellin transform by a polynomial corresponds to a differential operator in
\(x\,d/dx\)). Here ``even'' means \(g_0(1/x)=x\,g_0(x)\), equivalently \(B(1-z)=B(z)\). We fix \(B^{(M)}_w:=B^{(M)}\), the function of Lemma \ref{hyp:pw}, for every width \(w=2^{-j}\leq1\) of the grid; it satisfies \((\ast)\) at width \(w\) with \(C_M=(3/w)^M\), since \(1+t\geq w(1+t/w)\) for \(w\leq1\). Only the width \(w=1\) is used below. The
\emph{filled cell} is \(\tilde V_p:=V_p\oplus\C^{(D+1)^2}\)
with \(\tilde Q_p:=(Q_2|_{V_p})\oplus I\) (positive filler of dimension equal to
the cell's cap).
\end{definition}

\begin{definition}[localized scheme]\label{def:loc-scheme}
For each \(N\), the \emph{grid} \(\mathcal G_N\) is the finite computable set of
profiles \(p\) with \(D+3\leq M\leq N\), \(w,b\in\{2^{-j}:0\leq j\leq N\}\), and centres in
\((b\Z)^2\cap[-2^N,2^N]^2\). Set
\[
x_0(N):=\max_{p\in\mathcal G_N}\frac{\negind(\tilde Q_p)}{\dim\tilde V_p}
\in\mathbb Q\cap[0,\tfrac12],
\]
and \(\sigma_{\mathrm{loc}}:=\sst([x_0(N)])\). For each \(N\) fix a winning
profile \(p^*_N\) (maximizing the ratio), let \(\bar Q_N\) be the reduced
nondegenerate form of \(\tilde Q_{p^*_N}\) (quotient by the radical) and
\(\bar n_N:=\dim\tilde V_{p^*_N}-\dim\operatorname{rad}(\tilde Q_{p^*_N})\) its
dimension, and set
\[
\begin{gathered}
x_{\mathrm{base}}(N):=\frac{\negind(\bar Q_N)}{\bar n_N},\qquad
x_k(N):=\frac{\negind(\bar Q_N^{\otimes2^k})}{\bar n_N^{2^k}},\\[2pt]
\sigma_{\mathrm{base}}^{\,\mathcal U}:=\sst([x_{\mathrm{base}}(N)]),\qquad
\sigma_{\mathrm{loc}}^{(k)}:=\sst([x_k(N)]).
\end{gathered}
\]
Note that \(x_{\mathrm{base}}(N)\neq x_0(N)\) in general when the radical is
nontrivial (the denominators differ), and that
\(\sigma_{\mathrm{base}}^{\,\mathcal U}\) and the intermediate floors
\(\sigma_{\mathrm{loc}}^{(k)}\), \(k\geq1\), are not claimed to be independent
of \(\mathcal U\): only the limit dichotomy of Proposition
\ref{prop:loc-logistic} is.
\end{definition}

\begin{proposition}[well-definedness and autonomy]\label{prop:loc-welldef}
\textup{(a)} \(x_0(N)\) is a maximum of finitely many ratios, each determined by
a finite hermitian Gram matrix with entries \(Q_2(h,h')\), \(h,h'\) in a cell;
the cells are specified without zero positions. If the entries are expressed as
computable reals (which would follow from a prime-side expression of \(Q_2\),
not established here), the predicate \(\negind(\tilde Q_p)\geq1\) is
semidecidable: a strictly negative eigenvalue is certified at finite precision.
\textup{(b)} \(\sigma_{\mathrm{loc}}\in[0,\tfrac12]\) exists for every
\(\mathcal U\). \textup{(c)} \(x_0(N)\) is nondecreasing in \(N\)
(\(\mathcal G_N\subseteq\mathcal G_{N+1}\)), so it converges and
\(\sigma_{\mathrm{loc}}=\lim_N x_0(N)\) is independent of \(\mathcal U\).
\end{proposition}

\begin{proof}
(a) The grid is finite and each cell has dimension \(\leq(D+1)^2\); the filler
is an explicit block; a strictly negative eigenvalue of a hermitian matrix is
stable under perturbations of norm smaller than its modulus (Lemma
\ref{thm:weyl-stab}), so it is certified by a finite-precision approximation of
the entries. (b) \(\sst\) exists on limited elements. (c) Monotone bounded;
\(\sst\) of a convergent sequence is its limit, for every \(\mathcal U\).
\end{proof}

\begin{proposition}[the logistic law, localized]
\label{prop:loc-logistic}
For every \(\mathcal U\) and standard \(k\geq0\),
\[
\sigma_{\mathrm{loc}}^{(k)}
=\half\bigl(1-(1-2\sigma_{\mathrm{base}}^{\,\mathcal U})^{2^k}\bigr),
\qquad
\sigma_{\mathrm{base}}^{\,\mathcal U}\in
[\sigma_{\mathrm{loc}},\,2\sigma_{\mathrm{loc}}]\cap[0,\tfrac12].
\]
In particular, \emph{independently of \(\mathcal U\)}:
\(\lim_k\sigma_{\mathrm{loc}}^{(k)}=0\) if \(\sigma_{\mathrm{loc}}=0\), and
\(=\half\) if \(\sigma_{\mathrm{loc}}\in(0,\tfrac12]\).
\end{proposition}

\begin{proof}
For each \(N\), \(\bar Q_N\) is finite-dimensional nondegenerate; Theorem
\ref{thm:logistic} applies coordinatewise with base \(x_{\mathrm{base}}(N)\),
and \(\sst\) is a ring homomorphism on limited elements, giving the closed
formula in \(\sigma_{\mathrm{base}}^{\,\mathcal U}\). The bracketing: the filler
\(I_{(D+1)^2}\) is nondegenerate and orthogonal to the compressed block, so
\(\operatorname{rad}(\tilde Q_p)=\operatorname{rad}(Q_2|_{V_p})\oplus0\) and
\(\bar n_N\geq(D+1)^2\geq\tfrac12\dim\tilde V_{p^*_N}\), whence
\(\dim\tilde V_{p^*_N}/\bar n_N\in[1,2]\); since
\(\negind(\bar Q_N)=\negind(\tilde Q_{p^*_N})\) (the radical carries no negative
direction) and \(\sst\) is multiplicative on limited elements,
\(\sigma_{\mathrm{loc}}\leq\sigma_{\mathrm{base}}^{\,\mathcal U}
\leq2\sigma_{\mathrm{loc}}\). Moreover
\(x_{\mathrm{base}}(N)\leq(\bar n_N-(D+1)^2)/\bar n_N\leq\half\), so
\(\sigma_{\mathrm{base}}^{\,\mathcal U}\in[0,\half]\) and the case
\(\sigma^0=1\) (the only one breaking the dichotomy in Theorem
\ref{thm:logistic}) is excluded. The dichotomy: \(\sigma_{\mathrm{loc}}=0\iff
\sigma_{\mathrm{base}}^{\,\mathcal U}=0\) (bracketing), in which case the tower
is identically \(0\); and \(\sigma_{\mathrm{loc}}>0\Rightarrow
\sigma_{\mathrm{base}}^{\,\mathcal U}\in(0,\half]\Rightarrow
|1-2\sigma_{\mathrm{base}}^{\,\mathcal U}|<1\Rightarrow
\lim_k\sigma_{\mathrm{loc}}^{(k)}=\half\), for every \(\mathcal U\).
\end{proof}

\begin{remark}[normalization]\label{rem:radical-bug}
The closed formula of Proposition \ref{prop:loc-logistic} has
\(\sigma_{\mathrm{base}}^{\,\mathcal U}\), not \(\sigma_{\mathrm{loc}}\), as the
base of the tower. The two differ in general when the winning Gram is
degenerate, because \(x_0(N)\) is normalized by \(\dim\tilde V_p\) and
\(x_{\mathrm{base}}(N)\) by the reduced dimension \(\bar n_N\). The bracketing
\(\sigma_{\mathrm{base}}^{\,\mathcal U}\in[\sigma_{\mathrm{loc}},
2\sigma_{\mathrm{loc}}]\) suffices for the \(0\)-versus-\(\half\) dichotomy and
for the \(\mathcal U\)-independence of the limit; the intermediate floors are not
claimed to be canonical.
\end{remark}

\subsection{Stability of the cell index}
These three perturbation statements make a computable mesh suffice to capture a plane
whose exact position is unknown.

\begin{lemma}[Weyl stability {\cite{Bhatia}}]\label{thm:weyl-stab}
For hermitian \(A,E\) on a finite-dimensional space, with eigenvalues ordered
\(\lambda_1\geq\dots\geq\lambda_n\), \(|\lambda_k(A+E)-\lambda_k(A)|\leq\|E\|\)
for all \(k\). In particular, if \(A\) has \(\kappa\) eigenvalues \(\leq-\mu\)
and \(\|E\|<\mu\), then \(\negind(A+E)\geq\kappa\).
\end{lemma}

\begin{proof}
Courant--Fischer: \(\lambda_k(A+E)=\min_{\operatorname{codim}S=k-1}
\max_{0\neq v\in S}R_{A+E}(v)\) with \(R\) the Rayleigh quotient; for all \(v\),
\(|R_{A+E}(v)-R_A(v)|\leq\|E\|\). Take the optimal \(S\) for \(A\) and swap
\(A\leftrightarrow A+E\) \cite{Bhatia}.
\end{proof}

\begin{lemma}[bounded-rank stability; standard interlacing {\cite{Bhatia}}]\label{thm:rank-stab}
\(|\negind(A+E)-\negind(A)|\leq\rank E\).
\end{lemma}

\begin{proof}
Let \(N\) be negative for \(A\) with \(\dim N=\negind(A)\). On
\(N\cap\ker E\) (codimension \(\leq\rank E\) in \(N\)),
\((A+E)|_{N\cap\ker E}=A|_{N\cap\ker E}<0\), so
\(\negind(A+E)\geq\negind(A)-\rank E\). Swap \(A\leftrightarrow A+E\) and
\(E\leftrightarrow-E\).
\end{proof}

\begin{theorem}[Lipschitz Gram under cell shift]\label{thm:lipschitz}
For a profile \(p\), let \(G(\vartheta)\) be the Gram matrix of \(Q_2\) in the
basis \((h_{g_a,\varphi_e})\) of the cell with centres
\((\eta+\vartheta_1,\eta'+\vartheta_2)\). Then \(G\) is differentiable, and for
every compact \(K\subset\R^2\) there is \(L_{p,K}<\infty\) with
\(\|\partial_{\vartheta_i}G(\vartheta)\|\leq L_{p,K}\) for
\(\vartheta\in K\), \(i=1,2\).
\end{theorem}

\begin{proof}
Write \(B=B^{(M)}_w\), \(m=M-D\geq3\), and let \(|\vartheta|\leq R\) on \(K\).
By the Cauchy estimate on the unit circle about \(z\) and \((\ast)\),
\(|B'(z)|\leq C_Me^{a_0(|\operatorname{Re}z-1/2|+1)}(1+1/w)^M
(1+|\operatorname{Im}z|/w)^{-M}\), since
\(1+(|y|-1)_+/w\geq(1+|y|/w)/(1+1/w)\). Hence \(B\) and \(B'\) both satisfy
\((\ast)\) up to a constant depending on \(M,w,a_0\). On
\(|\operatorname{Re}\sigma|\leq1\), using \(|\sigma|^a\leq(1+|\operatorname{Im}
\sigma|)^D\cdot2^D\) and \(1+|t|\leq(1+|\eta|+R)(1+|t-\eta-\vartheta_1|)\),
each of \(\hat g_a\) and \(\partial_{\vartheta_1}\hat g_a=-iB'(\sigma-i(\eta+
\vartheta_1))\sigma^a\) is bounded by
\(C_K\max(1,w^{-1})^M(1+|\operatorname{Im}\sigma|)^{-m}\) uniformly for
\(\vartheta\in K\), and similarly on the \(\tau\) side. By the inequality
\eqref{eq:elem} (with \(\eta=\eta'=0\)), every term of \eqref{eq:speca} for a
Gram entry, and its \(\vartheta_i\)-derivative, is bounded on \(K\) by
\(C'_K(1+|\gamma_1|)^{-m}(1+|\gamma_2|)^{-m}\), which is summable over pairs of
zeros. By dominated convergence the entries are differentiable with derivatives
bounded on \(K\); the norm of a \((D+1)^2\times(D+1)^2\) matrix is at most its
dimension times its maximal entry.
\end{proof}

\begin{theorem}[the aligned cell survives in the mesh]\label{thm:mesh}
Let \(p^*\) be an ideal profile (real centres, not necessarily in the mesh) with
\(G(p^*)\) having \(\kappa\) eigenvalues \(\leq-\mu<0\), and let \(K\) be the
closed ball of radius \(2\) about \(0\) in the shift variable \(\vartheta\) of
Theorem \ref{thm:lipschitz}, applied at \(p^*\); put \(L:=L_{p^*,K}\). Let
\(p\in\mathcal G_N\) be the nearest mesh profile with the same \(D,w,M\) and
caliber \(b\leq1\), so \(|\mathrm{centres}(p)-\mathrm{centres}(p^*)|\leq
b\sqrt2\). If \(b\leq\mu/(4\sqrt2\,L)\), then \(\negind(G(p))\geq\kappa\).
In particular, for every ideal profile with \((\mu,L)\) fixed there is an
\(N_0\), explicit in \((\mu,L,\mathrm{position})\), such that \(\mathcal G_N\)
contains a cell of index \(\geq\kappa\) for all \(N\geq N_0\). \textup{(}We say
``explicit'', not ``computable'': the inputs \((\mu,L,\mathrm{position})\) are
not effectively knowable; cf.\ Section \ref{sec:honest}\textup{(a)}.\textup{)}
\end{theorem}

\begin{proof}
The segment from the centres of \(p^*\) to those of \(p\) lies in \(K\). By
Theorem \ref{thm:lipschitz} and the mean value theorem,
\(\|G(p)-G(p^*)\|\leq L(|\vartheta_1|+|\vartheta_2|)\leq L\cdot b\sqrt2\cdot2
=2\sqrt2\,Lb\leq\mu/2<\mu\); apply Lemma \ref{thm:weyl-stab}. Existence of
\(N_0\): the mesh has caliber \(2^{-N}\to0\) and reach \(2^N\to\infty\).
\end{proof}

\subsection{Realizability under \texorpdfstring{$\neg$}{not-}RH}
We now prove the realizability theorem for the doubled form \eqref{eq:speca}.
We use the zero-counting estimate \(N(t+1)-N(t)\ll\log(2+t)\)
\cite[Thm.~5.8]{IwaniecKowalski}.

\begin{lemma}[tail of the concentrated test]\label{lem:tail}
Let \(h=h_{g,\varphi}\) with \(|\hat g(\sigma)|\leq C_g(1+|\operatorname{Im}\sigma-
\eta|)^{-(M-D)}\) and \(|\hat\varphi(\tau)|\leq C_\varphi(1+|\operatorname{Im}
\tau-\eta'|)^{-(M-D)}\) on \(|\operatorname{Re}|\leq1\), for some
\(\eta,\eta'\in\R\). Let \(\Omega_r\) be the set of pairs with
\(\max(|\operatorname{Im}\sigma-\eta|,|\operatorname{Im}\tau-\eta'|)\geq r\).
If \(M-D\geq4\), then
\[
\sum_{(\rho_1,\rho_2)\in\Omega_r}|\hat h(\rho_1,\rho_2)|\,
|\hat h(\iota(\rho_1,\rho_2))|\leq
C_g^2C_\varphi^2\,\frac{C^*\,(\log(2+|\eta|+|\eta'|))^2}{(1+r)^{M-D-4}},
\]
with \(C^*\) an absolute constant. For an element of a cell,
\(\hat g=\sum_ax_a\hat g_a\) with centre \(\eta\) and coefficient vector of norm
\(\leq C_1\), and \(\hat\varphi=\sum_ey_e\hat\varphi_e\) with centre \(\eta'\)
and \(\|y\|\leq C_1\), one may take
\[
C_g=K_M(1+|\eta|)^D,\qquad C_\varphi=K_M(1+|\eta'|)^D,\qquad
K_M:=\sqrt{D+1}\,C_1C_Me^{3a_0/2}\max(1,w)^M .
\]
These constants depend on \(M\); for the bumps of Lemma \ref{hyp:pw} and
\(r\) above an absolute constant the right-hand side tends to \(0\) as
\(M\to\infty\) (Lemma \ref{hyp:pw}(iii)), while for the witness of Theorem \ref{thm:realiz} at radius \(2\) (\(r=\half\)) it diverges
for every family (Remark \ref{rem:pw-rigidity}).
\end{lemma}

\begin{proof}
Write \(u=\operatorname{Im}\sigma=(\gamma_1+\gamma_2)/2\),
\(v=\operatorname{Im}\tau=(\gamma_1-\gamma_2)/2\), so \(\gamma_1=u+v\),
\(\gamma_2=u-v\). With \(a=|\gamma_1-\eta-\eta'|\), \(b=|\gamma_2-\eta+\eta'|\),
one has \(a,b\leq|u-\eta|+|v-\eta'|\); the elementary inequality
\((1+\tfrac{a+b}2)^2\geq(1+a)(1+b)\) (since
\((1+\tfrac{a+b}2)^2-(1+a)(1+b)=(\tfrac{a-b}2)^2\geq0\)) then gives
\begin{equation}\label{eq:elem}
(1+|u-\eta|)(1+|v-\eta'|)\geq 1+|u-\eta|+|v-\eta'|\geq1+\tfrac{a+b}2
\geq\sqrt{(1+a)(1+b)}.
\end{equation}
The point \(\iota(\rho_1,\rho_2)\) has the same ordinates
(\(\operatorname{Im}(1-\bar\rho)=\operatorname{Im}\rho\)), so with \(m:=M-D\),
\[
|\hat h(P)||\hat h(\iota P)|\leq\tfrac14 C_g^2C_\varphi^2
\bigl[(1+|u-\eta|)(1+|v-\eta'|)\bigr]^{-2m}.
\]
On \(\Omega_r\) the product in brackets is \(\geq1+r\); extract
\((1+r)^{-(m-4)}\) and bound the remaining power \(-(m+4)\) by
\((1+a)^{-(m+4)/2}(1+b)^{-(m+4)/2}\) using \eqref{eq:elem}. The double sum over
pairs then factors into two single sums
\(\sum_\rho(1+|\gamma-c|)^{-k}\), \(k=(m+4)/2\geq4\), \(c=\eta\pm\eta'\);
grouped over unit intervals \(|\gamma-c|\in[j,j+1)\) with
\(N(t+1)-N(t)\ll\log(2+t)\) and
\(\log(2+|c|+j)\leq\log(2+|c|)\log(2+j)/\log2\), each is
\(\ll\log(2+|\eta|+|\eta'|)\sum_j(1+j)^{-k}\log(2+j)\), and the last series is
bounded uniformly in \(k\geq4\). This gives the bound with an absolute
\(C^*\). For the constants: on \(|\operatorname{Re}\sigma|\leq1\),
\(|\sigma^a|\leq(1+|\operatorname{Im}\sigma|)^D\leq(1+|\eta|)^D
(1+|\operatorname{Im}\sigma-\eta|)^D\) for \(a\leq D\); by \((\ast)\),
\(|B(\sigma-i\eta)|\leq C_Me^{3a_0/2}(1+|\operatorname{Im}\sigma-\eta|/w)^{-M}\);
\(1+t/w\geq(1+t)/\max(1,w)\); finally \(\sum_a|x_a|\leq\sqrt{D+1}\,\|x\|\). The
\(\tau\) side is the same with \(\eta'\) in place of \(\eta\).
\end{proof}

\begin{lemma}[an admissible family of bumps]\label{hyp:pw}
Fix an even \(\psi\in C_c^\infty(-1,1)\) with \(\psi\geq0\) and
\(\int\psi=1\), let \(\Psi(\zeta):=\int\psi(x)e^{-i\zeta x}\,dx\), and for
\(M\in\N\) put
\[
B^{(M)}(z):=\Bigl(\frac{\sin(\zeta/2)}{\zeta/2}\Bigr)^{M}\Psi(\zeta),
\qquad \zeta:=-i\bigl(z-\tfrac12\bigr).
\]
\textup{(i)} \(B^{(M)}\) is the Mellin transform of an even bump
\(g_0^{(M)}\in C_c^\infty(\R_{>0})\) \textup{(}even in the sense
\(g_0^{(M)}(1/x)=x\,g_0^{(M)}(x)\), equivalently \(B^{(M)}(1-z)=B^{(M)}(z)\)\textup{)}
with log-support in \([-a_0,a_0]\), and it satisfies \((\ast)\) with
\[
w=1,\qquad C_M=3^M,\qquad a_0=\tfrac M2+1 .
\]
\textup{(ii)} For every real \(x\), \(B^{(M)}(x)\) is real and
\(B^{(M)}(x)\geq1\).

\textup{(iii)} In the witness construction of Theorem \ref{thm:realiz}
\textup{(}radius \(R=4000\), width \(w=1\), centres
\((\eta^*,\eta'^*)=(\gamma_0,0)\), degree \(D\) fixed by the configuration and
\(R\)\textup{)}, the coefficient norm \(C_1\) is bounded independently of
\(M\), and the tail budget of Lemma \ref{lem:tail} at \(r=R/2\) satisfies
\[
C_g(M)^2\,C_\varphi(M)^2\;C^*\,(\log(2+\gamma_0))^2\;\Bigl(1+\frac R2\Bigr)^{-(M-D-4)}
\;\leq\;c\,\Bigl(\frac{81e^3}{2001}\Bigr)^{M}\;\longrightarrow\;0
\qquad(M\to\infty),
\]
with \(c\) depending on the configuration and \(R\) but not on \(M\), and
\(81e^3<81\cdot 20.1<1629\).
\end{lemma}

\begin{proof}
(i) Let \(\chi\) be the indicator function of \([-\half,\half]\) and
\(f:=\chi^{*M}*\psi\) (\(M\)-fold convolution). Then \(f\) is even, smooth (a
convolution with \(\psi\)), supported in \([-\tfrac M2-1,\tfrac M2+1]\), and
its Fourier transform \(\hat f(\zeta)=\int f(t)e^{-i\zeta t}dt\) is
\((\sin(\zeta/2)/(\zeta/2))^M\Psi(\zeta)\), since
\(\hat\chi(\zeta)=\sin(\zeta/2)/(\zeta/2)\). Put
\(g_0^{(M)}(x):=x^{-1/2}f(\log x)\in C_c^\infty(\R_{>0})\). Then
\(\int_0^\infty g_0^{(M)}(x)x^{z-1}dx=\int f(t)e^{t(z-1/2)}dt=\hat f(-\zeta)
=\hat f(\zeta)=B^{(M)}(z)\); evenness of \(f\) gives
\(g_0^{(M)}(1/x)=x\,g_0^{(M)}(x)\) and \(B^{(M)}(1-z)=B^{(M)}(z)\). For the
bound, note \(\operatorname{Re}\zeta=\operatorname{Im}z\) and
\(|\operatorname{Im}\zeta|=|\operatorname{Re}z-\half|\). From
\(\sin(\zeta/2)/(\zeta/2)=\half\int_{-1}^1e^{-i\zeta s/2}ds\) and
\(|\sin(\zeta/2)|\leq\cosh(\operatorname{Im}\zeta/2)\),
\[
\Bigl|\frac{\sin(\zeta/2)}{\zeta/2}\Bigr|\leq
e^{|\operatorname{Im}\zeta|/2}\min\Bigl(1,\frac2{|\operatorname{Re}\zeta|}\Bigr)
\leq\frac{3\,e^{|\operatorname{Re}z-1/2|/2}}{1+|\operatorname{Im}z|},
\]
because \(\min(1,2/t)\leq3/(1+t)\) for \(t\geq0\) (check \(t\leq2\) and
\(t\geq2\) separately). Also
\(|\Psi(\zeta)|\leq\int\psi(x)e^{|\operatorname{Im}\zeta||x|}dx\leq
e^{|\operatorname{Re}z-1/2|}\). Multiplying,
\(|B^{(M)}(z)|\leq3^Me^{(M/2+1)|\operatorname{Re}z-1/2|}(1+|\operatorname{Im}z|)^{-M}\),
which is \((\ast)\) with the stated constants.

(ii) For real \(x\), \(\zeta=iy\) with \(y=\half-x\) real, and
\(\sin(iy/2)/(iy/2)=\sinh(y/2)/(y/2)\geq1\) (value \(1\) at \(y=0\)), while
\(\Psi(iy)=\int\psi(t)\cosh(yt)\,dt\geq\int\psi=1\) since \(\psi\) is even and
nonnegative.

(iii) The witness of Theorem \ref{thm:realiz}, Step 2, is
\(\hat g_v(\sigma)=B^{(M)}(\sigma-i\gamma_0)\,P(\sigma)L(\sigma)\) and
\(\hat\varphi_v(\tau)=B^{(M)}(\tau)\,P'(\tau)/B^{(M)}(0)\), where \(P,P'\)
are polynomials determined by the configuration and \(L\) is the polynomial of
degree \(\leq1\) with \(L(\sigma_0)=1/(B^{(M)}(\half+\delta_0)P(\sigma_0))\)
and \(L(\sigma_0')=-1/(B^{(M)}(\half-\delta_0)P(\sigma_0'))\). By (ii) the three
values \(B^{(M)}(\half\pm\delta_0)\), \(B^{(M)}(0)\) are \(\geq1\), so the
coefficients of \(L\) are bounded by
\((|P(\sigma_0)|^{-1}+|P(\sigma_0')|^{-1})(1+(1+|\sigma_0|)/(2|\delta_0|))\),
and the coefficient vectors of \(PL\) and \(P'/B^{(M)}(0)\) in the bases
\((\sigma^a)\), \((\tau^e)\) have norm at most a constant \(C_1\) depending
on the configuration and \(R\) only. With \(w=1\), \(C_M=3^M\),
\(a_0=M/2+1\), Lemma \ref{lem:tail} gives
\(K_M=\sqrt{D+1}\,C_1e^{3/2}(3e^{3/4})^M\), \(C_g=K_M(1+\gamma_0)^D\),
\(C_\varphi=K_M\); hence the budget equals
\[
(D+1)^2C_1^4e^6(1+\gamma_0)^{2D}C^*(\log(2+\gamma_0))^2
\Bigl(1+\frac R2\Bigr)^{D+4}\cdot\Bigl(\frac{81e^3}{1+R/2}\Bigr)^{M},
\]
and \(1+R/2=2001>1629>81e^3\). The first factor does not depend on \(M\).
\end{proof}

\begin{remark}[the tail bound fails at radius \(2\)]\label{rem:pw-rigidity}
With the tail bound of Lemma~\ref{lem:tail}, the proof of Theorem~\ref{thm:realiz} does not go through at radius \(2\) (centres within \(\half\) of \((\gamma_0,0)\), budget at \(r=\half\), witness normalized by \(\hat\varphi_v(0)=1\)) for any family of bumps satisfying \((\ast)\), whatever \(w\), \(C_M\) and \(a_0(M)\). Indeed, write \(B=B^{(M)}_w\) and let the \(\tau\)-centre satisfy
\(|\eta'^*|<\half\). The witness has
\(\hat\varphi_v(0)=y_0B(-i\eta'^*)=1\), so \(C_1\geq|y_0|\geq
1/|B(-i\eta'^*)|\). Since \(\log|B|\) is subharmonic, \(\log|B(-i\eta'^*)|\)
is at most its mean over the unit circle about \(-i\eta'^*\), on which
\(\operatorname{Re}z=\cos\theta\) and
\(\operatorname{Im}z=\sin\theta-\eta'^*\). Inserting \((\ast)\) and
\(1+t/w\geq(1+t)/\max(1,w)\),
\[
\log\frac{C_M\max(1,w)^M}{|B(-i\eta'^*)|}\geq M\,A-a_0\cdot\tfrac32,
\qquad A:=\frac1{2\pi}\int_0^{2\pi}\log(1+|\sin\theta-\eta'^*|)\,d\theta .
\]
On the set \(|\sin\theta|\geq\tfrac45\), of proportion
\(1-\tfrac2\pi\arcsin\tfrac45>0.409\), one has
\(|\sin\theta-\eta'^*|\geq0.3\), so \(A\geq0.409\log1.3>0.107>\tfrac14\log\tfrac32\).
Hence \(K_M\geq C_1C_Me^{3a_0/2}\max(1,w)^M\geq e^{AM}\), and the budget at
\(r=\half\) is at least a positive constant times
\(K_M^4(3/2)^{-M}\geq(e^{4A}/\tfrac32)^M\to\infty\). (Numerically the
minimum of \(A\) over \(|\eta'^*|\leq\half\) is \(0.473\ldots\), attained at
\(\eta'^*=0\), so the budget then grows at least like \((4.4)^M\), and the
same argument, with centres within \(\half\) of \((\gamma_0,0)\), shows that
the budget at radius \(R\) can tend to \(0\) only if
\(1+R/2>e^{4\cdot0.473}\approx6.6\).) This concerns the upper bound of Lemma
\ref{lem:tail} only; it is why Theorem \ref{thm:realiz} works with a radius
above an absolute constant, and \(R=4000\) is a convenient choice, not an
optimal one.
\end{remark}

\begin{lemma}[Vandermonde with bump]\label{lem:vandermonde}
Let \(z_1,\dots,z_q\in\C\) be distinct and \(B\) entire with \(B(z_i)\neq0\). If
\(D\geq q-1\), the evaluations \(f\mapsto f(z_i)\) are linearly independent on
\(\operatorname{span}\{B(\sigma)\sigma^a:0\leq a\leq D\}\); for any data
\((y_1,\dots,y_q)\) there is an element with \(f(z_i)=y_i\) and coefficient norm
\(\leq C(z,\min_i|B(z_i)|)\|y\|\). Moreover, for
\(B=B^{(M)}_w(\cdot-i\eta)\not\equiv0\), the centres \(\eta\) with
\(B(z_i-i\eta)=0\) for some \(i\) are discrete; every open interval of centres
contains an \(\eta\) with \(\min_i|B(z_i-i\eta)|\geq b_0>0\), stable on a ball.
\end{lemma}

\begin{proof}
The matrix \((B(z_i)z_i^a)_{i,a}=\operatorname{diag}(B(z_i))\cdot(z_i^a)\) has
rank \(q\) (Vandermonde with distinct nodes, invertible diagonal). The function
\(\eta\mapsto B(z_i-i\eta)\) is entire and not identically zero, so its real zeros are
discrete; a finite union of discrete sets is discrete; \(B\) is continuous and
nonzero on a compact set avoiding the zeros.
\end{proof}

\begin{theorem}[realizability under \(\neg\)RH]
\label{thm:realiz}
Assume \(\neg\)RH. Let
\(\rho^0=\half+\delta_0+i\gamma_0\) be an off-critical
zero, with \(\gamma_0>0\), \(\delta_0\neq0\) and multiplicity \(m_0\geq1\). Then there are a standard
dimension \(D_0\) and a standard \(N_0\) such that for all \(N\geq N_0\),
\(x_0(N)\geq 1/(2D_0)>0\), and hence for every non-principal \(\mathcal U\),
\[
\sigma_{\mathrm{loc}}\geq\frac{1}{2D_0}>0,\qquad
\lim_{k\to\infty}\sigma_{\mathrm{loc}}^{(k)}=\tfrac12.
\]
\end{theorem}

\begin{proof}
\emph{Step 0 (target plane).} Take \(P_0=(\rho^0,\rho^0)\) with coordinates
\((\sigma_0,\tau_0)=(\half+\delta_0+i\gamma_0,0)\) and
\(\iota P_0=(1-\bar\rho^0,1-\bar\rho^0)\) with
\((\sigma_0',0)=(\half-\delta_0+i\gamma_0,0)\). Since \(\delta_0\neq0\),
\(\iota P_0\neq P_0\): a free orbit, contributing
\(2m_0^2\operatorname{Re}[\hat h(P_0)\overline{\hat h(\iota P_0)}]\) to
\eqref{eq:speca}.

\emph{Step 1 (the sets to annihilate).} Let \(S\) be the finite set of zeros
\(\rho\) with \(||\gamma|-\gamma_0|\leq R\), where \(R:=4000\), and \(Z_\sigma\) the finite set of
\(\sigma=(\rho_1+\rho_2)/2\) over pairs with both members in \(S\), excluding
\(\sigma_0,\sigma_0'\). Let \(Z_\tau\) be the finite set of
\(\tau=(\rho_1-\rho_2)/2\) over the \emph{resonant pairs}: pairs with both
members in \(S\) and \(\sigma\in\{\sigma_0,\sigma_0'\}\), other than the target
orbit \(\{P_0,\iota P_0\}\). Every resonant pair has \(\tau\neq0\) (\(\tau=0\)
together with \(\sigma=\sigma_0\) forces \(\rho_1=\rho_2=\rho^0\), the target).
Put \(D:=\max(|Z_\sigma|+1,\,|Z_\tau|)\); all standard, fixed by the
configuration and \(R\), and independent of \(M\). \textup{(}Under general \(\neg\)RH configurations the resonant
pairs contribute terms of uncontrolled sign, which is why they are
annihilated.\textup{)}

\emph{Step 2 (the cell and witness).} Choose the ideal profile \(p^*\) with
degree \(D\), width \(w=1\) (so \(B=B^{(M)}\) of Lemma \ref{hyp:pw}), decay
order \(M\geq D+4\), \(\sigma\)-centre \(\eta^*=\gamma_0\) and \(\tau\)-centre
\(\eta'^*=0\). Let \(P(\sigma):=\prod_{z\in Z_\sigma}(\sigma-z)\) and
\(P'(\tau):=\prod_{z\in Z_\tau}(1-\tau/z)\) (recall \(0\notin Z_\tau\)), and let
\(L\) be the polynomial of degree \(\leq1\) with
\(L(\sigma_0)=1/(B(\half+\delta_0)P(\sigma_0))\) and
\(L(\sigma_0')=-1/(B(\half-\delta_0)P(\sigma_0'))\); it exists because
\(\sigma_0\neq\sigma_0'\), \(P(\sigma_0)P(\sigma_0')\neq0\)
(\(\sigma_0,\sigma_0'\notin Z_\sigma\)) and \(B(\half\pm\delta_0)\geq1\)
(Lemma \ref{hyp:pw}(ii)). In \(V_{p^*}\) take \(v=h_{g_v,\varphi_v}\) with
\[
\hat g_v(\sigma):=B(\sigma-i\gamma_0)P(\sigma)L(\sigma),\qquad
\hat\varphi_v(\tau):=B(\tau)P'(\tau)/B(0);
\]
these lie in the cell since \(\deg(PL)\leq|Z_\sigma|+1\leq D\) and
\(\deg P'=|Z_\tau|\leq D\), and \(B(0)\geq1\). Since
\(\sigma_0-i\gamma_0=\half+\delta_0\) and \(\sigma_0'-i\gamma_0=\half-\delta_0\),
we have \(\hat g_v=0\) on \(Z_\sigma\), \(\hat g_v(\sigma_0)=1\),
\(\hat g_v(\sigma_0')=-1\), \(\hat\varphi_v=0\) on \(Z_\tau\) and
\(\hat\varphi_v(0)=1\); no nonvanishing of \(B\) on \(Z_\sigma\cup Z_\tau\) is
needed. The coefficient norm \(C_1\) is independent of \(M\) (Lemma
\ref{hyp:pw}(iii)), and \(\hat g_v,\hat\varphi_v\in\mathcal E\) because
\(M\geq D+3\). The data \((1,-1)\) fixes the
phase: the target plane contributes
\(2m_0^2\operatorname{Re}[(\half\cdot1\cdot1)\overline{(\half\cdot(-1)\cdot1)}]
=-m_0^2/2<0\).

\emph{Step 3 (budget of \(Q_2(v,v)\)).} Decompose \eqref{eq:speca}: (a) the orbit
\(\{P_0,\iota P_0\}\) contributes \(-m_0^2/2\); (b) pairs with both members in
\(S\) other than the target: either \(\sigma\in Z_\sigma\), where
\(\hat g_v=0\), or \(\sigma\in\{\sigma_0,\sigma_0'\}\) (resonant), where
\(\tau\in Z_\tau\) and \(\hat\varphi_v=0\) -- in both cases the contribution is
\(0\) (the case split is exhaustive, and the \(\iota\)-image adds no new case:
\(\sigma(\iota P)=1-\bar\sigma\) swaps \(\sigma_0\leftrightarrow\sigma_0'\),
already covered); (c) the remaining pairs have a member \(\rho_i\) outside \(S\), so
\(|\gamma_i-\gamma_0|\geq||\gamma_i|-\gamma_0|>R\); since \(\gamma_{1,2}=u\pm v\), this gives
\(|u-\eta^*|+|v-\eta'^*|=|u-\gamma_0|+|v|>R\), and the pair lies in
\(\Omega_{R/2}\) about \((\eta^*,\eta'^*)\). By Lemma \ref{lem:tail} with
\(r=R/2\) and the constants stated there, their total is at most the budget
of Lemma \ref{hyp:pw}(iii), which tends to \(0\) as \(M\to\infty\) (the
zero-counting estimate enters only through \(C^*\)). Choose \(M=M_0\geq D+4\) standard
with \(|(c)|\leq m_0^2/4\). Then
\[
Q_2(v,v)\leq-\frac{m_0^2}{2}+0+\frac{m_0^2}{4}=-\frac{m_0^2}{4}<0,
\]
so \(\negind(Q_2|_{V_{p^*}})\geq1\), with a standard spectral margin \(-\mu_0\).

\emph{Step 4 (ideal cell to grid).} By Theorem \ref{thm:mesh} (with
\(\kappa=1\), \(\mu=\mu_0\), \(L=L_{p^*,K}\) finite by Theorem \ref{thm:lipschitz},
all standard) there is a standard \(N_0\) such that for \(N\geq N_0\),
\(\mathcal G_N\) contains a profile \(p\) with \(\negind(Q_2|_{V_p})\geq1\).

\emph{Step 5 (conclusion).} For \(N\geq N_0\),
\(x_0(N)\geq\negind(\tilde Q_p)/\dim\tilde V_p\geq1/(2(D+1)^2)=:1/(2D_0)\) with
\(D_0=(D+1)^2\) standard (filler doubles dimension; the negative subspace
persists).
By Proposition \ref{prop:loc-welldef}(c),
\(\sigma_{\mathrm{loc}}=\lim x_0(N)\geq1/(2D_0)\) for every \(\mathcal U\);
Proposition \ref{prop:loc-logistic} gives \(\sigma_{\mathrm{loc}}^{(k)}\to\half\).
\end{proof}

\begin{remark}[the conceptual key]
\label{rem:key}
The index \(\negind(Q_2|_{V_p})\) is a property of the \emph{compressed Gram}.
Test-accessibility restricts the \emph{specification} of \(V_p\) (no zero
positions); the \emph{witness} \(v\) inside \(V_p\) is an object of the proof,
free to use all positions. Polynomial factors with roots at zero positions are
obtained as linear combinations \emph{inside} the test-accessible span
\(\{B\cdot\sigma^a\}\) (Lemma \ref{lem:vandermonde}). Localizing to the bounded
height of the first off-critical zero removes the regime \(T\to\infty\), but
the proof still needs a localization radius above an absolute constant and
bumps whose constants grow only exponentially in \(M\) (Lemma \ref{hyp:pw},
Remark \ref{rem:pw-rigidity}). We claim no quantitative bound: \(D_0,N_0\)
are not effective -- they depend on the unknown local configuration.
\end{remark}

\begin{proposition}[RH branch]\label{prop:RHbranch}
Under RH, \(x_0(N)=0\) for all \(N\), \(\sigma_{\mathrm{loc}}=0\), and
\(\sigma_{\mathrm{loc}}^{(k)}\equiv0\).
\end{proposition}

\begin{proof}
Under RH every pair of zeros is fixed by \(\iota\), so
\(Q_2(h,h)=\sum_P|\hat h(P)|^2\geq0\) on \(\mathcal H\) (see the definition of
\(Q_2\) before Definition \ref{def:cell}); every compression of a positive form is positive; the
filler is positive; \(\negind=0\) in every cell.
\end{proof}

\begin{theorem}[the localized logistic dichotomy]\label{thm:main}
For every non-principal ultrafilter \(\mathcal U\):
\begin{enumerate}
\item[\textup{(i)}]
\(\mathrm{RH}\Longrightarrow\sigma_{\mathrm{loc}}=0\Longrightarrow
\lim_k\sigma_{\mathrm{loc}}^{(k)}=0\);
\item[\textup{(ii)}] \(\neg\mathrm{RH}\Longrightarrow\sigma_{\mathrm{loc}}>0\Longrightarrow
\lim_k\sigma_{\mathrm{loc}}^{(k)}=\tfrac12\).
\end{enumerate}
Hence
\[
\mathrm{RH}\;\Longleftrightarrow\;\sigma_{\mathrm{loc}}=0\;\Longleftrightarrow\;
\lim_{k\to\infty}\sigma_{\mathrm{loc}}^{(k)}=0\;\Longleftrightarrow\;
\lim_{k}\sigma_{\mathrm{loc}}^{(k)}\neq\tfrac12 .
\]
Moreover, \(\lim_k\sigma_{\mathrm{loc}}^{(k)}\in\{0,\half\}\) with
\(0\iff\sigma_{\mathrm{loc}}=0\), independently of \(\mathcal U\).
\end{theorem}

\begin{proof}
(i) Proposition \ref{prop:RHbranch} and Proposition \ref{prop:loc-logistic}.
(ii) Theorem \ref{thm:realiz} and Proposition \ref{prop:loc-logistic}. The
\(\mathcal U\)-independence of \(\sigma_{\mathrm{loc}}\) is Proposition
\ref{prop:loc-welldef}(c); the unconditional dichotomy is Proposition
\ref{prop:loc-logistic}.
\end{proof}

\begin{remark}[independence of \(\mathcal U\)]
The dichotomy of Theorem \ref{thm:main} is independent of the ultrafilter:
\(\sigma_{\mathrm{loc}}\) is a genuine limit (Proposition
\ref{prop:loc-welldef}(c)). The non-canonicity one might fear -- ``which
ultrafilter concentrates where the defect lives?'' -- is resolved by
monotonicity of the growing grid, not by a choice.
\end{remark}

\section{Scope and open problems}\label{sec:honest}

\subsection{(a) RH is not proved}
Theorem \ref{thm:main} is an equivalence between RH and
\(\sigma_{\mathrm{loc}}=0\) for the form \(Q_2\) of \eqref{eq:speca}, a
reformulation of RH. Because \(Q_2\) is defined by a sum over the zeros, it is
close to an elementary restatement of ``there is an off-critical zero'': its
only content is that such a zero is detected by a cell specified without zero
positions, with the witness and the tail estimate of Section
\ref{sec:logistic}. It does not decide RH and provides no new route to
deciding it. The standard part
\(\sigma_{\mathrm{loc}}=0\) is a \(\Pi_1\) statement over the sequence
\((x_0(N))_N\) -- ``for all large \(N\), \(x_0(N)\) small'' -- as non-effective
as RH itself. Each \(x_0(N)\) is finite linear algebra on Gram matrices of
\(Q_2\) (Proposition \ref{prop:loc-welldef}(a)). If the entries were expressed
through primes, a negative eigenvalue certified by interval arithmetic would
falsify RH; that \(\neg\)RH produces such an eigenvalue is Theorem
\ref{thm:realiz}.

\subsection{(b) Computing the symbol is a uniform statement over windows}
The trichotomy (Section \ref{sec:trichotomy}) locates the obstruction at the
single point \(0\in\spec\sigma(Q)\), and the following identity shows that
deciding whether the symbol equals its autonomous prediction amounts to a
uniform second-order asymptotic over windows, which is not known.

\begin{theorem}[computing the symbol \(=\) uniformity over windows]
\label{thm:symbolcompute}
On the same field of test spaces,
\(\|\sigma(Q)-\sigma(\Qaut)\|_{\Cor}=\limsup_n\|Q_n-\Qaut_n\|\). Hence the
statement ``the symbol of the Weil form is the archimedean/near-prime
prediction'' is equivalent to \(\|Q_n-\Qaut_n\|\to0\): the second-order
deficit asymptotics, in operator norm per window, \emph{uniformly} over \(n\) on
the whole ball of tests.
\end{theorem}

\begin{proof}
Lemma \ref{lem:symbol-norm} applied to the difference section; the right-hand
side is the operator norm of the per-window deficit form.
\end{proof}

The right-hand side is a uniform second-order pair-correlation asymptotic
over windows, which is not known. Mean-value methods address at most the
\emph{Ces\`aro states} (the mean over windows) of \(\sigma(Q)\); strict
positivity is a statement of \emph{norm} (all states, including the pure states
supported on window subsequences aligned with the zeros). In this language,
the missing step is the passage from mean states to all states of the corona.
We record the abstract half of what is
provable.

\begin{proposition}[mean Ces\`aro positivity, abstract half]\label{prop:cesaro}
If, for every \(\varepsilon>0\),
\(N^{-1}\#\{n\leq N:\lambda_{\min}(Q_n)<-\varepsilon\}\to0\) (``almost every
window is \(\varepsilon\)-positive''), then for every positive test section
\(0\leq A_n\leq I\) and every Banach limit \(\mathcal L\),
\[
\mathcal L\text{-}\lim_N\frac1N\sum_{n\leq N}\frac{\tr(Q_nA_n)}{d_n}\geq0:
\]
all positive Ces\`aro states of \(\sigma(Q)\) are \(\geq0\).
\end{proposition}

\begin{proof}
If \(\lambda_{\min}(Q_n)\geq-\varepsilon\), then \(Q_n\geq-\varepsilon I\), so
\(\tr(Q_nA_n)\geq-\varepsilon\tr(A_n)\geq-\varepsilon d_n\) (using
\(\tr(CA)\geq0\) for \(C,A\geq0\)); otherwise the trivial bound
\(\tr(Q_nA_n)\geq-\sup_k\|Q_k\|\,d_n\). The bad windows have vanishing density,
so the Banach limit of a sequence \(\geq-\varepsilon-o(1)\) is
\(\geq-\varepsilon\); let \(\varepsilon\to0\).
\end{proof}

The hypothesis is of mean-value / large-sieve type (almost every window); it
is not proved here for the Weil windows. Proposition \ref{prop:cesaro}
guarantees that once supplied, it yields ``mean positivity of the symbol''.
In this language, the missing step is the passage from mean states to all
states.

\subsection{(c) What is new and what is standard}
\begin{itemize}
\item The window algebra \(\W\) is the \(\ell^\infty\)-product of the matrix
algebras \(\mathcal B(H_n)\) and \(\Cor\) its corona, a standard construction;
what we use is its reading of the uniform-versus-mean quantifier (Corollary
\ref{cor:master-quantifier}).
\item Theorem \ref{thm:exact} is standard; we include a proof for completeness.
\item The trichotomy (Theorems \ref{thm:fat}--\ref{thm:nogo}) inside the model
\(\mathfrak W\) (Remark \ref{rem:trichotomy}); part (a) of Theorem
\ref{thm:nogo} uses the additional axiom (V3).
\item The logistic identity (Theorem \ref{thm:logistic}), its localized form
(Proposition \ref{prop:loc-logistic}), and Theorem \ref{thm:main}, whose RH
and \(\neg\)RH branches are both proved for the zero-side form \(Q_2\), the
latter with the explicit bumps of Lemma \ref{hyp:pw}.
\end{itemize}

\subsection{What remains open}
The computability of the symbol -- uniform pair-correlation asymptotics (Theorem
\ref{thm:symbolcompute}) -- is not addressed. Theorem \ref{thm:main} is
proved for the zero-side form \(Q_2\) and is close to a restatement of the
existence of an off-critical zero; it would become an arithmetic statement only
through a prime-side expression of \(Q_2\).
Conjecture \ref{conj:lp134} (resolution repulsion) would make strict symbol
positivity at one calendar equivalent to \(m<\infty\) in the model; with
replication (R), weak positivity would give \(m=0\). A prime-side expression of \(Q_2\) is not
established. We list the dependencies.

\begin{remark}[dependencies]\label{rem:deps}
\textup{Theorem \ref{thm:main}} depends on: the definition \eqref{eq:speca} of
\(Q_2\) and its absolute convergence (proved in Section \ref{sec:logistic});
the zero-counting estimate \(N(t+1)-N(t)\ll\log(2+t)\)
\cite[Thm.~5.8]{IwaniecKowalski}; and, for the \(\neg\)RH branch only,
Lemma \ref{hyp:pw}, proved in Section \ref{sec:logistic}. \textup{The trichotomy (Section
\ref{sec:trichotomy})} is proved inside the model \(\mathfrak W\); its transfer
to \(\zeta\) requires (S1)--(S3), (V1)--(V2) (and (V3) where used) and
Hypotheses \ref{hyp:blocks}--\ref{hyp:autonomy} for the Weil windows, none of
which is proved. \textup{Proposition \ref{prop:cesaro}} is unconditional given
its stated density hypothesis.
\end{remark}

\appendix
\section{The constructive regime: Weil positivity for finite systems of primes}
\label{app:constructive}

The Weil--Toeplitz analysis of the body concerns the \emph{defect} --- the off-line
zeros and the negative directions they create. It is illuminating to record, by way
of contrast, the opposite extreme: the fully constructive regime, in which positivity
is immediate and the whole content lies in a single decidable discriminator. The
material of this appendix is classical (the easy direction of Weil's criterion
\cite{Weil1952} for finite Euler products); we claim no novelty for it.

\subsection{Pure Euler objects}
For a prime $p$ write $L_p=\log p$. The \emph{Euler object} of $p$ with parameter
$\alpha$ ($1\le|\alpha|\le p$) is the datum $F_{p,\alpha}$ whose local factor is
$L(s)=(1-\alpha p^{-s})^{-1}(1-(p/\bar\alpha)p^{-s})^{-1}$, so that
$-L'(s)/L(s)=L_p\sum_{m\ge1}c_mp^{-ms}$ with $c_m=\alpha^m+(p/\bar\alpha)^m$.
Put $c_0:=2$ and $c_{-m}:=\overline{c_m}$. The poles of $L$ lie on
$\Re s=\log|\alpha|/L_p$ and on $\Re s=1-\log|\alpha|/L_p$, a configuration
invariant under $s\mapsto1-\bar s$. The admissible class is $C_c^\infty(\R)$
(complex valued); for $f$ admissible put $\tilde f(x)=\overline{f(-x)}$,
$(f\star\tilde f)(y)=\int f(x)\overline{f(x-y)}\,dx$ and
$\widehat f(\xi)=\int f(x)e^{-ix\xi}dx$. For $h\in C_c^\infty(\R)$ define the
(finite) sum
\[
W_{F_{p,\alpha}}(h):=\sum_{m\in\Z}c_m\,p^{-|m|/2}\,h(mL_p),
\]
and say that $F_{p,\alpha}$ satisfies the positivity property $H$ if
$Q_{F_{p,\alpha}}(f,f):=W_{F_{p,\alpha}}(f\star\tilde f)\ge0$ for all admissible
$f$. For a finite set of primes with Euler objects $X_p$, the product $X$ has
logarithmic derivative the sum of those of its factors, and we set
$W_X:=\sum_pW_{X_p}$.

\begin{theorem}[Local purity is the exact discriminator]\label{thm:local-purity}
For $1\le|\alpha|\le p$,
\[
F_{p,\alpha}\text{ satisfies }H\quad\Longleftrightarrow\quad|\alpha|=\sqrt p
\quad\Longleftrightarrow\quad|c_m|\le2p^{m/2}\ \ \forall m\ge1 ,
\]
and these hold if and only if every pole of the local factor lies on
$\Re s=\half$.
\end{theorem}
\begin{proof}
The last equivalence is the location of the poles. If $|\alpha|=\sqrt p$, write
$\alpha=\sqrt p\,e^{i\theta}$; then $p/\bar\alpha=\alpha$, $c_m=2\alpha^m$,
$|c_m|=2p^{m/2}$, and $c_mp^{-|m|/2}=2e^{im\theta}$ for all $m\in\Z$. With
$k=f\star\tilde f\in C_c^\infty(\R)$, $\widehat k=|\widehat f|^2$, the Poisson
summation formula gives
\[
Q_{F_{p,\alpha}}(f,f)=2\sum_{m\in\Z}e^{im\theta}k(mL_p)
=\frac{2}{L_p}\sum_{\xi\in\Lambda_p}|\widehat f(\xi)|^2\ge0,\qquad
\Lambda_p=\{(2\pi\ell-\theta)/L_p:\ell\in\Z\}.
\]
Conversely, let $|\alpha|=p^{1/2+\delta'}$ with $\delta'\neq0$. Then
$|c_m|\ge p^{m/2}|p^{m\delta'}-p^{-m\delta'}|>2p^{m/2}$ for large $m$, and
$c_mp^{-|m|/2}=e^{im\theta}\,2\cosh(m\delta'L_p)$ for $m\in\Z$. Fix $n\ge2$ and
a real $u\in C_c^\infty(\R)$ with $\|u\|_2=1$ supported in $(-\varepsilon,\varepsilon)$,
$\varepsilon<L_p/4$, and put $f=\sum_{j=0}^{n-1}t_ju(x-jL_p)$. Then
$k(mL_p)=\sum_{j-l=m}t_j\bar t_l$, and
\[
Q_{F_{p,\alpha}}(f,f)=2\sum_{j,l}s_j\bar s_l\cosh((j-l)\delta'L_p),\qquad
s_j=t_je^{ij\theta}.
\]
The matrix $C=(\cosh((j-l)x))_{j,l<n}$, $x=\delta'L_p$, equals
$\half(ab^{T}+ba^{T})$ with $a_j=e^{jx}$, $b_j=e^{-jx}$; its smallest eigenvalue
is $\half(n-|a||b|)\le\half(n-e^{(n-1)|x|})$, which is negative already for
$n=2$ (it equals $1-\cosh x$). So $H$ fails.
\end{proof}

\subsection{Additivity}
\begin{proposition}[Convolution does not mix Weil data]\label{prop:crosszero}
For Euler objects $X_1,X_2$ of distinct primes, the Weil functional of the product is
additive, $W_{X_1\star X_2}=W_{X_1}+W_{X_2}$; there are no cross terms.
Consequently a product of Euler objects over a finite set of primes satisfies
$H$ if and only if each factor is pure \textnormal{(}$|\alpha_p|=\sqrt p$\textnormal{)}.
\end{proposition}
\begin{proof}
The logarithmic derivative of a product is the sum of the logarithmic
derivatives and $W$ is linear in it, so $W_{X_1\star X_2}=W_{X_1}+W_{X_2}$; a sum
of positive forms is positive, which gives ``if'' by Theorem
\ref{thm:local-purity}. For ``only if'', let the factor at $q$ be impure,
$|\alpha_q|=q^{1/2+\delta'}$, $\delta'\neq0$, and let $r$ be the number of other
primes. Take $f=\sum_{j<n}t_ju(x-jL_q)$ as in the proof of Theorem
\ref{thm:local-purity}, with $\varepsilon$ so small that
$|mL_{p}-\ell L_q|>2\varepsilon$ for all other primes $p$ of the system, all
$m\neq0$ with $|m|L_p\le nL_q+1$ and all $|\ell|<n$ (possible since $p^m\neq q^\ell$).
Then $k=f\star\tilde f$ vanishes at every $mL_p$, $m\neq0$, and
$k(0)=\sum_j|t_j|^2$, so each other prime contributes $2\sum_j|t_j|^2$, and
\[
Q_X(f,f)=2r\|s\|^2+2\,s^*Cs .
\]
Choosing $s$ an eigenvector of $C$ for its smallest eigenvalue
$\le\half(n-e^{(n-1)|\delta'|L_q})$ makes this negative for $n$ large.
\end{proof}

\subsection{Scope}
Proposition~\ref{prop:crosszero} with Theorem~\ref{thm:local-purity} is the
\emph{easy} direction of Weil's criterion \cite{Weil1952} in a finite model:
purity places every pole on $\Re s=\half$ \emph{by hypothesis}, and positivity
follows by Poisson summation, with the additivity of logarithmic derivatives
making a finite system no harder than one prime. We claim no novelty for this
material. Purity is the exact discriminator: $H$ holds if and only if
$|\alpha_p|=\sqrt p$ for every $p$ in the system. The finite Euler product is
\emph{not} $\zeta$, and the contrast with the body is the point: in this regime
purity is assumed and positivity is immediate, whereas the main text concerns the
absence of an a priori purity.

\begin{thebibliography}{99}

\bibitem{Bhatia}
R.~Bhatia,
\emph{Matrix Analysis},
Graduate Texts in Mathematics \textbf{169}, Springer, 1997.

\bibitem{BottcherSilbermann}
A.~B\"ottcher and B.~Silbermann,
\emph{Analysis of Toeplitz Operators},
2nd ed., Springer, 2006.

\bibitem{Gantmacher}
F.~R.~Gantmacher,
\emph{The Theory of Matrices}, Vol.~2,
Chelsea, 1959.

\bibitem{IwaniecKowalski}
H.~Iwaniec and E.~Kowalski,
\emph{Analytic Number Theory},
American Mathematical Society Colloquium Publications \textbf{53}, AMS, 2004.

\bibitem{Kato}
T.~Kato,
\emph{Perturbation Theory for Linear Operators},
2nd ed., Springer, 1976.

\bibitem{RabinovichRochSilbermann}
V.~Rabinovich, S.~Roch, and B.~Silbermann,
\emph{Limit Operators and Their Applications in Operator Theory},
Operator Theory: Advances and Applications \textbf{150}, Birkh\"auser, 2004.

\bibitem{ReedSimon}
M.~Reed and B.~Simon,
\emph{Methods of Modern Mathematical Physics, IV: Analysis of Operators},
Academic Press, 1978.

\bibitem{Robinson}
A.~Robinson,
\emph{Non-standard Analysis},
North-Holland, 1966.

\bibitem{Weil1952}
A.~Weil,
\emph{Sur les ``formules explicites'' de la th\'eorie des nombres premiers},
Comm.\ S\'em.\ Math.\ Univ.\ Lund, Tome suppl\'ementaire (1952), 252--265.

\bibitem{P18} D.~A.~Trejo Pizzo, \emph{The localized Weil form on even test functions: off-line Lorentzian blocks and obstructions to positive-kernel realizations}, preprint, 2026 (Paper 18 of this series).

\bibitem{P25} D.~A.~Trejo Pizzo, \emph{The Weil form and the scaling operator: a Pontryagin-index dictionary in zero coordinates}, preprint, 2026 (Paper 25 of this series).

\end{thebibliography}

\end{document}
